\chapter{Slopes, Tangents, and Rates of Change}
\label{sec:veloc-slop-tang}
% \baselineskip=1.5cm
% \section*{The Principle of Local Linearity}
% \label{sec:LocalLinearity}
% \begin{wrapfigure}[10]{r}{2.2in}
%   \vskip-4mm{}
%   \captionsetup{labelformat=empty}
%   \centerline{\includegraphics*[height=1.6in,width=2.2in]{../Figures/DifferentialTriangle1}}
%   \label{fig:DiffTri}
% \end{wrapfigure}
% So far we have used our differentiation rules to turn an ordinary
% equation like $y=x^2,$ into the differential equation
% $\dx{y}=2x\dx{x}.$ All well and good but of what use is this? What, if
% anything, does the differential equation tell us about the formula $y=x^2?$

% As you are surely aware when you look at a very small section of any
% curve it is almost indistinguishable from a straight line. Moreover,
% the smaller the section the more it resembles a straight line segment.
% This fact is known as the \underline{\term{Principle of Local
%     Linearity}} and is illustrated by the sketch at the right. The red
% curve is the graph of $y=x^2$ ``blown up'' near the point $(a,a^2)$ so
% that we only see a very small part of it. As you can plainly see the
% hypotenuse of the blue triangle is
% \begin{enumerate}
% \item A straight line and, 
% \item Practically coincident with the
%   curve.
% \end{enumerate}

\section{Slopes and Tangents}
\label{sec:slopes-tangents}
\markright{\sc Slopes and  Tangents}


As we observed in Chapter~\ref{sec:some-rules-diff} Newton, Leibniz,
and their contemporaries extended the Principle of Local Linearity to
its logical limit. They said if you cut out an infinitely small
section of a curve then the section actually is an infinitely short
straight line segment. Of course we are quite familiar with infinitely
short line segments. We call them differentials.

\begin{wrapfigure}[12]{r}{3in}
  \vskip-6mm{}
  \captionsetup{labelformat=empty}
  \centerline{\includegraphics*[height=2.2in,width=3in]{../Figures/DifferentialTriangle1}}
  \label{fig:DiffTri}
\end{wrapfigure}
Leibniz called a triangle whose sides are differentials, like the blue
triangle in the figure at the right, a \term{differential triangle}\aside{Because,
  what else would you call it?}.  In the figure $\dx{y}$ and
$\dx{x}$ are differentials in the vertical and horizontal directions,
respectively. But $\dx{s}$ is a differential at the point $(a, a^2)$ in the direction of the
curve  or, what comes to  the same thing, in the
direction tangent to the curve. Thus $\dx{s}$ is a differential
along the curve, and simultaneously  along the line
tangent to the curve\aside{The idea that $\dx{s}$ is a
  straight line segment on the curve and also on its tangent
    line is a bit  troubling. It is one of many bizarre consequences
  of  the notion of differentials.}.

% To generate a formula for this tangent line is straightforward. We
% already have a point on the tangent line: $(a,a^2)$. All we need is
% the slope of $\dx{s}$ at $(a,a^2)$ and we can use the Point-Slope
% Formula.


The differentials $\dx{x}$, and $\dx{y}$, don't really tell us much
about the graph of $y=x^2$, but the \term{differential ratio},
$\dfdx{y}{x}$, does. At each point, $(a, a^2)$, $\dfdx{y}{x}$ gives us
the slope of $\d{s}$ and the line tangent to the curve at that point.
Loosely speaking, the differential ratio $\dfdx{y}{x}$ gives the slope
of the curve at each point.

% slope of $\d{s}$ at a given point will be the same
% as the slope of the curve at the same point.  It should be clear that the slope of the
% line tangent to a curve will change from point to point, unless the
% curve is itself a straight line.  Thus thee \term{differential ratio}
% $\dfdx{y}{x}$, tells us the slope of the graph at each point of the
% graph.

If we want to find the slope of the curve $y=x^2$ at the point
$(a,a^2)$ we simply find $\d{y}$  as usual:$ \dx{y}=2x\dx{x}$.  Then
we take the extra step of dividing through by $\dx{x}$ to get
$\dfdx{y}{x}=2x$.  Evaluating this at the point $x=a$, $y=a^2$ will give us
the slope of the line tangent to the graph of $y=x^2$ at the point
$(a,a^2)$.

\begin{ProficiencyDrill}
  Find an equation of the line tangent to the graph of $y=x^2$ at
  each of the following points.
  \begin{enumerate}[label={\bf{}(\alph*)}]
    \begin{multicols}{5}
    \item $(-2,4)$
    \item $(-1,1)$
    \item $(0,0)$
    \item $(1,1)$
    \item $(2,4)$
    \end{multicols}
  \end{enumerate}
    Compare these results to those obtained using Fermat's Method in
  Problem~\#\ref{prob:ferm-meth-x2}
\end{ProficiencyDrill}

\begin{mynotation}{Evaluation}
  We will frequently need to evaluate the expression $\dfdx{y}{x}$ at
  different points. This can  quickly become  very confusing
  unless we have some way of indicating which point is under
  consideration. To avoid  confusion, we  use the notation \\
\centerline{
$\eval{\dfdx{y}{x}}{(x,y)}{(a,b)}$ or $\geneval{\dfdx{y}{x}}{(a,b)}$.
}
This notation is very flexible. Frequently the $y$ coordinate will not
be in play. In that case we write \\
\centerline{$\eval{\dfdx{y}{x}}{x}{a}$}  
  to indicate that we are evaluating $\dfdx{y}{x}$ at the
  point $(x,y)=(a,b)$.
\end{mynotation}
\begin{myexample}
    For example to compute the slope of the graph of $y=3x^2-5x$ at
    $(3,12)$ we use the following three-step process.
    \begin{enumerate}
    \item First  differentiate $y=3x^2-5x$ to obtain the
      differential equation $\dx{y}=(6x-5)\dx{x}$. This 
      relates $\dx{y}$ and $\dx{x}$.
    \item Second, from this differential equation we 
      find the ratio  $\dfdx{y}{x}=6x-5.$ This
      differential ratio tells us the slope of the curve at
      every point on the curve.
    \item Third, if we need the value of $\dfdx{y}{x}$ at a single
      point like $(3,12)$, for example\aside{Obviously this makes no
        sense if the point $(3,12)$ is not a point on the curve. You
        should always make sure the problem makes sense before you
        try to solve it.}, we compute
      $$\eval{\dfdx{y}{x}}{(x,y)}{(3,12)}=\bigeval{6
        x-5}{(x,y)}{(3,12)}=6\cdot3-5=13.$$

   In this case we
      could just write: $\eval{\dfdx{y}{x}}{x}{3}=13$, since the $y$
      coordinate never comes into play.
    \end{enumerate}
    Finally, we emphasize that Leibniz' notation is
    \emph{deliberately} evocative of the notion of slope because when
    we evaluate the differential ratio at a point it tells us the
    slope of the curve at that point.  In view of the \alert{Principle
      of Local Linearity}, this is equivalent to finding the slope of
    the line tangent to the curve at that point.
  \end{myexample}

  This notation probably seems unnecessarily cumbersome and honestly, in
  some ways it is. But be assured we have not made this choice
  lightly, and this notation has certain advantages that will become
  apparent later.

  It can be easy to get careless or simply confused and write $\dfdx{y}{x}$
  when you mean $\eval{\dfdx{y}{x}}{x}{a}$, and \foreign{vice versa},
  so it is important that you have the distinction between them clear
  in your mind.

\begin{embeddedproblem-enumerate}
  \begin{enumerate}[label={\bf{}(\alph*)}]
    \item Explain the difference between $\dfdx{y}{x}$ and
      $\eval{\dfdx{y}{x}}{x}{a}$ carefully and clearly.
    \item We (the authors) have sometimes had students assert that
      since $\dfdx{y}{x}=2x$, the equation of the line tangent to the
      graph of $y=x^2$ at the point $(a, a^2)$ is
      \begin{equation}
      y-a^2=2x(x-a).\label{eq:TanLineNoLine}
    \end{equation}
    Set $a=3$. Is this the equation of a line?
  \item Find the correct formula for line tangent to the graph of
    $y=x^2$ at the point $(a, a^2)$.  
  \end{enumerate}
%  \vskip-5mm{}
\end{embeddedproblem-enumerate}
%\end{mynotation}

\begin{ProficiencyDrill}
  Find $\eval{\dfdx{y}{x}}{(x,y)}{(2,2)}$ if $y=\sqrt{2x}$ and compare this with the slope
  we obtained in problem~\#\ref{prob:DescartesNormals2} using
  Descartes' Method of Normals.
\end{ProficiencyDrill}


  \begin{embeddedproblem}
    \label{embedded:Deri}
    Evaluate $\eval{\dfdx{y}{x}}{x}{a}$ for each of the following
    functions at each of the given values of $x$.
    \begin{enumerate}[label={\bf{}(\alph*)}]{}
    \item $y=x^3$, $a=0,\pm1,\pm2$ (Compare with
      Problem~\#\ref{prob:ferm-meth-x2}.)
    \item $y=x^3+2x$, $a=0,\pm1,\pm2$ (Compare with
      Problem~\#\ref{prob:ferm-meth-aedeq-1}.)
    \item $y=\sqrt{x}$, $a=4$ (Compare with
      Problem~\#\ref{prob:DescartesNormals1}.)
    \end{enumerate}
    % This is why the title for Leibniz' paper was \foreign{``{\bf A
    % New
    % Method for} .\ .\ . {\bf{}Tangents} Which is Impeded
    % Neither by Fractional nor by Irrational Quantities, and a
    % Remarkable Type of Calculus for This.''}
  \end{embeddedproblem}

%\pagebreak[4]

There is real temptation at this point to cut corners and compute the
differential ratio
$\dfdx{y}{x}$ all in one step. On purely practical grounds, we urge you to
differentiate first and then divide by $\dx{x}$ in two separate
steps. Our reasons are two-fold.  First, the rules we've learned are
\emph{differentiation} rules. They were designed to compute
differentials, not differential ratios.  Second, in applications
to come we will see that dividing by $\dx{x}$ will not always move the
solution process forward. It can sometimes even get in the way,
depending on the problem.

Nonetheless, as you gain more experience and a deeper understanding
you will find yourself gravitating toward the differential ratio more
and more. This is a normal progression. As we have said before
differentials are an intermediate step, a convenient fiction. You
don't want to become tied to them for life.

% \subsection{Defining the Tangent Line}
\section{Defining the Tangent Line}
\label{sec:defin-tang-line}
\begin{wrapfigure}[9]{l}{1in}
  \vskip-6mm{}
  \captionsetup{labelformat=empty}
  \centerline{\includegraphics*[height=1.5in,width=1in]{../Figures/TangentLineDef6}}
  \label{fig:TangentLineDef6}
\end{wrapfigure}
The word tangent comes from the Latin word \foreign{tangere} (to
touch).  So you'd think it would be easy to define what we mean by
``tangent line.'' When asked, most people will come up with something
like, ``The line tangent to a curve is the line that touches the curve
at exactly one point.'' Does this seem like a reasonable definition to
you? Stop and think about it for a few minutes. We're obviously
offering a plausible sounding definition that we know won't work. Do
you see the problem?

It certainly \emph{seems} OK, and if we check it with a few simple
examples everything looks reasonable. 
But consider a small segment of the graph of $y=x^3$
at the point $(1,1)$ and its corresponding tangent line as seen at
the left.

\begin{wrapfigure}[]{r}{1in}
  \vskip-3mm{}
  \captionsetup{labelformat=empty}
  \centerline{\includegraphics*[height=1.7in,width=1in]{../Figures/TangentLineDef7}}
  \label{fig:TangentLineDef7}
\end{wrapfigure}
While this certainly fits our notion of a tangent line touching a
curve only once, look what happens if we zoom out (see the sketch at
the right).  This supposed tangent line will cross (touch) the
graph of the curve a second time.  In fact, there is only one point on
the graph of $y=x^3$ where the tangent line touches the graph only
once. Do you see which point that is?  Give this a few minutes
thought.

In view of what we've just seen, what do you think the definition of a
tangent line
\emph{ought} to be?

In the figure at the right it is clear that the ``touching'' is
different at the two points. At $(1,1)$ the blue tangent line
``grazes'' the graph of $y=x^3$ at $(1,1)$ exactly as we would
expect. But at the other point it actually cuts through the graph.

 The difficulty seems to be that we said ``touch'' when we meant
``graze''.  How about this definition instead: ``The line tangent to a
curve is the line that \emph{grazes} the curve at exactly one point.''

If it is not possible for a line to ``graze'' a curve at more than one
point this could be a workable definition. Unfortunately the next
drill shows that this is not true.

%\pagebreak[4]
\begin{ProficiencyDrill}
  \begin{wrapfigure}[]{r}{2.5in}
  \vskip-8mm{}
    \captionsetup{labelformat=empty}
    \centerline{\includegraphics*[height=1in,width=2in]{../Figures/TangentLineDef4}}
    \label{fig:TangentLineDef5}
  \end{wrapfigure}
  Show that $ y=\frac12x+\frac14$ is the equation of the
  line tangent to the graph of
  $$
  y=\frac12x+\frac14 + (x-1)^2(x-2)^2(x-3)^2(x-4)^2.
  $$
  at   $x=1, 2, 3, $ and $x=4$  as seen in the
  sketch at the right.% \\
  % \centerline{\includegraphics*[height=2in,width=3in]{../Figures/TangentLineDef4}}
\end{ProficiencyDrill}


It appears that clearly defining what we mean by ``tangent line'' is
not as simple as it seemed to be at first. This is frustrating because
we all have a pretty clear intuitive idea of what we mean. The difficulty seems to be capturing in
words the mental image we all share.

\begin{wrapfigure}[5]{left}{1.in}
  \vskip-4mm{}
  \captionsetup{labelformat=empty}
  \centerline{\includegraphics*[height=.8in,width=1.2in]{../Figures/DifferentialTriangle3}}
  \label{fig:DiffTri3}
\end{wrapfigure}
Recall that we have informally stated   the \term{Principle of Local
  Linearity} as: ``Locally any curve is indistinguishable from a straight line.''

But which straight line?

The answer to that question is visible in the differential triangle at
the left. It is clear that, \emph{locally} (near a given point $(x,y)$
on the curve) a curve will look like the hypotenuse of the
differential triangle $\d{s}$ at the point $(x,y)$.  So we make the
following definition.

\begin{mydefinition}[The Line Tangent to a Curve, at a Point]
  \label{def:TangentLine}
  The \term{line tangent to the graph of a function} $y(x)$, at the point
  $(a,y(a))$ is the line which passes through the point $(a,y(a))$
  with slope $m=\eval{\dfdx{y}{x}}{x}{a}$.
\end{mydefinition}

Do you  see what we did here? Intuitively we know what properties we want a
tangent line to have. We want it to ``touch'' or ``graze'' our
curve at a particular point but we've seen that neither of those words
quite captures everything we need. And what we need is for the tangent
line to pass through the tangent point and have the same slope as the
function at that point. So we've removed the words ``touch'' and
``graze'' in favor of the more precise, quantifiable language of Calculus. 

% aren't precise enough. So we simply
% defined the phrase ``tangent line'' in terms of the properties we
% wanted. But ``touch'' and ``graze'' aren't precise enough so we
% removed them. The language we chose tied the tangent idea to a
% particular point; we made it local.
% By our definition the blue line
% above is tangent to the graph of $f(x)=x\sin(x)$ \emph{at the point}
% $(7.96,f(7.96))$ (approximately) because it passes through
% $(7.96,f(7.96))$ and it's slope is equal to
% $\eval{\dfdx{(x\sin(x))}{x}}{x}{7.96}$.

Since we \emph{defined} the tangent line to be the extension of
$\dx{s}$ to the finite scale we have essentially defined it to have
exactly that property we will need to make Calculus useful.  Moreover
with this definition in place the \term{Principle of Local Linearity}
can be formalized.

\begin{mydefinition}[The Principle of Local Linearity]
  \label{def:LocalLinearity}
  Sufficiently close to a given point every smooth curve is
  indistinguishable from the line tangent at that point.
\end{mydefinition}

The history of mathematics is riddled with this sort of ``backward
looking'' definition and we will see it again. For example Newton and
Leibniz both knew what they wanted their \foreign{Calculus
  Differentialus} (Leibniz) and \emph{Method of Fluxions} (Newton) to
do. They worked out their computational procedures intuitively -- just
as we started with an intuitive understanding of ``tangent line.'' But
coming up with precise definitions and a solid logical foundation was
harder. It took about $200$ years for the mathematical community to
devise definitions that gave back what intuition had led us to. This
is part of why we are not yet overly concerned about our inability to
precisely define a differential. When we finally address this issue in
Chapter~\ref{chapt:whats-wrong-with} we will not be devising any new
differentiation techniques. We will be providing logical underpinnings
for the techniques our intuition has led us to and that our experience
has shown to be useful.

\pagebreak[4]
\begin{ProficiencyDrill}
  Find an equation of the line tangent to the given curve at the
  indicated values of $x$.
  \begin{enumerate}[label={\bf{}(\alph*)}]
    \begin{multicols}{3}
    \item $ y =\frac{2x}{x+1}$\\ when $ x =1.$
    \item $y=\dfrac{1}{1+x^2}$\\ when $ x=-1$
    \item $y=x+\sqrt{x}$\\ when $x=1$
    \item $y^2=x^3$\\ when $x=1$
    \item $y^3=x^2$\\ when $x=1$
    \item $y=\dfrac{\sqrt{x}}{x+1}$\\ when $x=4$
    \item $y = 3x$\\ when $x=10$
    \item $y = -7x$\\ when $x=-100$
    \item $y = -\dfrac{x^2}{\pi}$\\ when $x=\sqrt{\pi}$
    \item $y=(3x-1)^{-6}$\\ when $x=0$
    \item
      $y = x^3-\dfrac1x$\\ when $x=1$
    \item $y=\sqrt{1+x^2}$\\ when $x=2$
    \item $y = \dfrac{x}{x-2}$\\ when $x=0$
    \item $y = \dfrac{x}{x-2}$\\ when $x=\frac14$
    \item $y = \dfrac{1}{1+x^2}$\\ when $x=\frac12$
    \item $y = \dfrac{1}{1-x^2}$\\ when $x=1$
    \item $y = \dfrac{1}{\sqrt{x}}$\\ when $x=9$
    \item $y = \dfrac{1}{\sqrt[3]{x^2}}$\\ when $x=4$
    \item $y=\sqrt{2x+\sqrt[3]{64-x}}$\\ when $x=0$
    \item $y = \pi x^2$\\ when $x=\dfrac{1}{\sqrt{\pi}}$
    \end{multicols}
  \end{enumerate}
%  \vskip-4mm{}
\end{ProficiencyDrill}

  \begin{embeddedproblem}
    \label{problem:tanlinedef}
    Show that the only point on the graph of $y=x^3$ where the tangent
    line touches the curve \emph{only} at the point of tangency is the
    point $(0, 0)$ and that the tangent is horizontal.

    Now graph the curve and the line tangent at $(0,0)$ on the same
    set of axes. Does this \emph{look} like a tangent line to you?
  \end{embeddedproblem}
  
Does the graph help you see the significance of Problem~\#\ref{problem:tanlinedef}? The
line tangent to the graph of $y=x^3$ at the point $(0,0)$
actually crosses the graph at that point.
This is weird. And it is exactly counter to our original understanding
of what it means for a line to be tangent to a curve isn't it?

The next problem will probably strike you as equally weird. But in
both problems the weirdness is a logical consequence of our definition
of tangent line, so we will accept it as long as it doesn't lead to
inconsistencies.

Weirdness is acceptable. It is inconsistency that we need to avoid.

%\pagebreak[4]
\begin{embeddedproblem}
%  \vskip-8mm{}
  Use Definition~\ref{def:TangentLine} to find the equation of the
  line tangent to the graph of the straight line $y=mx+b$ at the point
  $(a, y(a))$. Show that this is in fact, the original line.
    % Suppose that the graph of the equation $y=mx+b$ passes through the
    % point $(x_0, y_0)$.
    % \begin{enumerate}[label={\bf{}(\alph*)}]
    %   % \begin{multicols}{2}
    % \item Re-express $y=mx+b$ in Point-Slope form using the point  $(x_0, y_0)$.
    % \item Find the equation of the line tangent to the graph of
    %   $y=mx+b$ at $(x_0, y_0)$ in Point-Slope form.
    % \item What does this tell you about the relationship between the
    %   graph of $y=mx+b$ and its tangent lines?
    %   % \end{multicols}
    % \end{enumerate}
  \vskip-6mm{}
\end{embeddedproblem}
%\vskip-4mm{}

\begin{embeddedproblem-enumerate}
  \begin{enumerate}[label={\bf{}(\alph*)}]
  \item Find the equations of all lines tangent to the graph of
    $4x^2+4y^2=25$ that are also parallel to the line $2x-3y=7$.\\
    \hint{There are two such lines, but you should find four possible
      candidates. You'll need to eliminate two of them.}
  \item   Find the equations of all lines tangent to the graph of
    $16x^2+9y^2=144$ that are also parallel to the line $8x+6y=8.$
  \item   Find the equations of all lines tangent to the graph of
    $16x^2-9y^2=144$ that are also parallel to the line $8x+6y=8.$
  \end{enumerate}
%  \vskip-3mm{}
\end{embeddedproblem-enumerate}


\begin{embeddedproblem}
  Find those points on the graph of the curve $y=x^3+5$  where the
  tangent line is
  \begin{enumerate}[label={\bf{}(\alph*)}]
    % \begin{multicols}{2}
  \item Parallel to the line: $12x-y=17$.
  \item Perpendicular to the line: $x+3y=2$.
    % \end{multicols}
  \end{enumerate}
  \vskip-6mm{}
\end{embeddedproblem}

\vskip1mm{}
\begin{embeddedproblem}[Find the Pattern]
  Find all points where the
  line tangent to the curve passes through the indicated points. If no
  such points exist, explain how you know.
  \begin{enumerate}[label={\bf{}(\alph*)}]
  \item    $y=x^2+2$  through the points
    \begin{enumerate}[label={\bf{}(\roman*)}]
      \begin{multicols}{4}
      \item $(0,0)$
      \item $(-5,0)$
      \item $(2,6)$
      \item $(1,6)$
      \end{multicols}
    \end{enumerate}
  \item    $y=\frac{x+1}{x-1}$  through the points
    \begin{enumerate}[label={\bf{}(\roman*)}]
      \begin{multicols}{4}
      \item $(0,0)$
      \item $(0,-5)$
      \item $(5,5)$
      \item $(-1,0)$
      \end{multicols}
    \end{enumerate}
  \item $y=\frac{x^2-1}{x+2}$  through the points
    \begin{enumerate}[label={\bf{}(\roman*)}]
      \begin{multicols}{4}
      \item $(0,0)$
      \item $(0, -1/2)$
      \item $(-5,-10)$
      \item $(-5,10)$
      \end{multicols}
    \end{enumerate}
  \end{enumerate}
  \vskip-4mm{}
\end{embeddedproblem}


\pagebreak[4]
\begin{embeddedproblem}
%\vskip-8mm{}
  \label{problem:ClassicalCurves1}
  Find an equation of the line(s) tangent to and normal to each of the
  given curves at the specified points.
  \begin{enumerate}[label={\bf{}(\alph*)}]{}
    \begin{multicols}{2}
    \item{\bf{}\href{https://mathshistory.st-andrews.ac.uk/Curves/Astroid/}{Astroid}:}\\ $x^\frac23+y^\frac23 = 4$ \\at
      $\left(1,3\sqrt{3}\right)$ and
      $\left(-1,-3\sqrt{3}\right).$\\
      \leftline{\includegraphics*[height=1.5in,width=1.5in]{../Figures/Astroid}}
    \item{\bf{}An Elliptic Curve:}\\
      $y^2=x^3-x+1$ \\at $(0,1)$ and $(0,-1)$.\\
      \leftline{\includegraphics*[height=1.5in,width=1.5in]{../Figures/EllipticCurve1}}
      % Note: This and more general elliptic curves (over finite fields)
      % form the basis of Elliptic Curve Cryptography and are useful in
      % many areas of number theory.
    \item{\bf{}\href{https://mathshistory.st-andrews.ac.uk/Curves/Cardioid/}{Cardioid}:}\\ $\left(x^2+y^2\right)^2+4x(x^2+y^2)=4y^2$ \\at
      $(0,2)$ and $(0,-2)$.\\
      \leftline{\includegraphics*[height=1.5in,width=1.5in]{../Figures/Cardioid}}
    \item{\bf{}\href{https://mathshistory.st-andrews.ac.uk/Curves/Conchoid/}{Conchoid}% \aside{\href{https://mathshistory.st-andrews.ac.uk/Biographies/Nicomedes/}{Nicomedes} (circa 280 BC - 210 BC) used
          % this curve to trisect an angle, a construction which is
          % impossible to accomplish with a straight edge and compass
          % alone.}
        of
        \href{https://mathshistory.st-andrews.ac.uk/Biographies/Nicomedes/}{Nicomedes}:}\
      \\ $(y-2)^2(x^2+y^2)=2y^2$ \\at $(1,1)$
      and $(3,3).$\\
      \leftline{\includegraphics*[height=1.5in,width=2in]{../Figures/ConchoidOfNicomedes}}
\item{\bf{}\href{https://mathshistory.st-andrews.ac.uk/Biographies/Cassini/}{Cassini}% \aside{\href{https://mathshistory.st-andrews.ac.uk/Biographies/Cassini/}{Giovanni Cassini} (1625-1712) studied Cassini Ovals
          % in astronomy.  Whereas ellipses are points the sum of whose
          % distances from two fixed foci is constant, Cassini ovals
          % are where the product of the distances is constant.  In
          % honor of his accomplishments in astronomy, the Cassini
          % space probe was launched in 1997 and became the first probe
          % to orbit Saturn. }
        \href{https://mathshistory.st-andrews.ac.uk/Curves/Cassinian/}{oval}:}\\
      $(x^2+y^2+1)^2-4x^2=\frac32$\\
      at
      $\left(\frac{\sqrt[4]{2}}{2},\frac{\sqrt[4]{2}}{2}\right)$\\
      \leftline{\includegraphics*[height=1.2in,width=2in]{../Figures/CassiniOval}}
    \item{\bf{}The \href{https://mathshistory.st-andrews.ac.uk/Curves/Bicorn/}{Bicorn}:}\\ $y^2\left(1-x^2\right)=\left(x^2+2y-1\right)^2$\\
      at
      $\left(\frac12, \frac{12\pm3\sqrt[]{3}}{26}\right)$\\
      \leftline{\includegraphics*[height=1.5in,width=1.in]{../Figures/Bicorn}}
    \end{multicols}
  \end{enumerate}
\end{embeddedproblem}

%\pagebreak[4]
\addtocounter{problem}{-1}
\begin{embeddedproblem-enumerate}[(continued)]
  \label{problem:ClassicalCurves2}
  % Find an equation of the line(s) tangent to and normal to each of the
  % given curves at the given points. (This is a continuation of the
  % previous problem.)
  \begin{enumerate}[label={\bf{}(\alph*)}]
    \setcounter{enumi}{6}
    \begin{multicols}{2}
% \item{\bf{}\href{https://mathshistory.st-andrews.ac.uk/Biographies/Cassini/}{Cassini}% \aside{\href{https://mathshistory.st-andrews.ac.uk/Biographies/Cassini/}{Giovanni Cassini} (1625-1712) studied Cassini Ovals
%           % in astronomy.  Whereas ellipses are points the sum of whose
%           % distances from two fixed foci is constant, Cassini ovals
%           % are where the product of the distances is constant.  In
%           % honor of his accomplishments in astronomy, the Cassini
%           % space probe was launched in 1997 and became the first probe
%           % to orbit Saturn. }
%         \href{https://mathshistory.st-andrews.ac.uk/Curves/Cassinian/}{oval}:}\\
%       $(x^2+y^2+1)^2-4x^2=\frac32$\\
%       at
%       $\left(\frac{\sqrt[4]{2}}{2},\frac{\sqrt[4]{2}}{2}\right)$\\
%       \leftline{\includegraphics*[height=1.2in,width=2in]{../Figures/CassiniOval}}
%     \item{\bf{}The \href{https://mathshistory.st-andrews.ac.uk/Curves/Bicorn/}{Bicorn}:}\\ $y^2\left(1-x^2\right)=\left(x^2+2y-1\right)^2$\\
%       at
%       $\left(\frac12, \frac{12\pm3\sqrt[]{3}}{26}\right)$\\
%       \leftline{\includegraphics*[height=1.5in,width=1.in]{../Figures/Bicorn}}
    \item{\bf{}\href{https://mathshistory.st-andrews.ac.uk/Curves/Kampyle/}{The
          Kampyle} of \href{https://mathshistory.st-andrews.ac.uk/Biographies/Eudoxus/}{Eudoxus}:}\\ $x^4=x^2+y^2$\\
      at
      $\left(\pm 2, \pm 2\sqrt{3}\right)$\\
      \leftline{\includegraphics*[height=1.5in,width=1.5in]{../Figures/KampyleEudoxus}}
    \item{\bf{}The
        \href{https://mathshistory.st-andrews.ac.uk/Curves/Folium/}{Folium}:}\\
      $(x^2+y^2)(y^2+x(x+4))=4xy^2$\\
      at
      $\left(-1, -1\right)$ and $(-1, 1)$\\
      \leftline{\includegraphics*[height=1.5in,width=2in]{../Figures/Folium}}
    \end{multicols}
  \end{enumerate}
\end{embeddedproblem-enumerate}


\begin{Digression}[Dividing by Zero]
% In general curves are very complex and we need to be careful with
% them.
\begin{ProficiencyDrill-1line}
  \label{drill:ZeroDenominator1}
  Show that if  $y^2=x^3-3x+1$ then $\dfdx{y}{x}=\frac{3x^2-3}{2y}$.
\end{ProficiencyDrill-1line}

%\pagebreak[4]
\begin{wrapfigure}[12]{r}{2in}
  \vskip-4mm{}
\captionsetup{labelformat=empty}
\centerline{\includegraphics*[height=2in,width=2in]{../Figures/EllipticCurve2}}
\label{fig:EllipticCurve2}
\end{wrapfigure}
In Drill~\#\ref{drill:ZeroDenominator1} it is tempting to to say that
$\dfdx{y}{x}=\frac{3x^2-3}{2y}$ will give us the slope of the tangent
line tangent to the curve at \emph{any} point $(x,y).$ Typically this would
be correct.  But not always.

For example  the tangent line of this curve  is horizontal
(has slope equal to zero) when
$$
0=\dfdx{y}{x}=\frac{3x^2-3}{2y}.
$$
Solving this gives $x=\pm1$. It is plainly
visible on the graph at the right that there is  a horizontal tangent line
when $x=-1$.
But notice that there is no point on the graph where $x=1$! 
Obviously it is meaningless to ask for the slope of a curve at a point that
isn't on the curve.

Even if we didn't have the graph to look at we could still tell that
at $x=1$ there is no tangent line.  Setting $x=1$ and trying to solve for
the corresponding $y$ values we get $y^2=1-3+1=-1$. Since there is no
real number that satisfies the equation $y^2=-1$ we conclude that
there is no point on the graph with $x$-coordinate equal to $1$.

On the other hand, there are three values of $x$ that correspond to
$y=0$.  These would be difficult to compute but from the graph we can
see that the tangent line is most likely vertical at all three points.
To compute $\dfdx{y}{x}=\frac{3x^2-3}{2y}$ at each of these points
we would have to divide by $y=0$.

Obviously when we divide a number by
zero we get infinity, and just as obviously the slope of a vertical
line is also infinity so everything fits. Right?
No, of course not. This is just a very deceptive coincidence.

Many
students come into a Calculus with the impression that
dividing by zero results in either zero or infinity.
  It does not. Division by zero is not meaningful.

  Mathematicians say that division by zero is an \term{undefined
    operation} or simply that it is \term{undefined}. At first it is
  not at all clear why this is, so let's take look at this
  question.

  What does it mean to divide by any number? To take a very concrete
  example, what does it mean to divide a length of five
  units\aside{The units could be inches, picometers, or
    light-years. We don't care.} by seventeen? It means we want to
  divide the length into seventeen equally sized pieces, right?  But
  if we want to divided $5$ units by zero, obviously we can't divide
  five into zero pieces. So in this very concrete example the phrase
  ``five divided by zero'' is simply and utterly without meaning.

  From a slightly more abstract viewpoint we could ask for the length
  of each piece when we divide a length of five into seventeen pieces.
  We get $\frac{5}{17}$ of course. That's what the fraction notation
  means\aside{Just as  $\frac12$ is the number we get when we
  divide one into two pieces.}. But what number do we get when we divide
  five into zero pieces. Once again the idea has no meaning.

  Taking an even more abstract point of view,  observe that multiplication and division are inverse
  operations (they undo each other). Twelve divided by three is four
  $\left(\frac{12}{3}=4\right)$ \emph{precisely because} when we
  multiply three times four we get twelve ($3\times 4 =12$).

  Now don't think of five as a length to be divided, think of it
  instead as simply a number. So what number can we multiply
  by zero to get five? Obviously there isn't any such number so this
  is a meaningless question.

  No matter what we try or how we look at it we cannot find any
  way to make  division by zero meaningful.
  

  \centerline{\alert{Division by zero is meaningless.}}
  \centerline{\alert{Don't ever try to divide by zero.}}
  \centerline{\alert{You will only irritate your
      teacher.}}
\end{Digression}

Since we can't divide by zero we will need to be a bit more careful
about how we interpret the symbol $\dfdx{y}{x}$. It will give us the
slope of the curve (equivalently, the slope of the tangent line) as
long as $\dfdx{y}{x}$ has meaning. But when $y$ is zero
$\dfdx{y}{x}=\frac{3x^2-3}{2y}$ is meaningless.  Thus it appears that
for this particular curve we can draw no conclusions whatsoever from
our formula about the tangent line when $y=0$. In this case we say
that the differential quotient is \term{undefined}. In general, the
differential quotient is \term{undefined} at any point where
evaluating $\dfdx{y}{x}$ requires a division by zero.

Does this mean we can draw no conclusions at all? Certainly not. In
fact those values of $x$ or $y$ where $\dfdx{y}{x}$ is undefined will
often turn out to be the most useful. But we are not yet quite prepared to
see how to deal with them properly so we will defer this discussion
until a later point.

At the very least this warns us that there are subtleties that must be
dealt with when using a powerful tool like Calculus. We will avoid
these subtleties for a bit longer, but you need to know that they
exist. Don't try to force matters by, for example, blindly dividing by
zero as if it is meaningful. We will return to these matters when we
have a better understanding of the underlying principles.

\begin{embeddedproblem}
% \begin{wrapfigure}[]{r}{.75in}
%   \captionsetup{labelformat=empty}
%   \vskip-13mm{}
%   \centerline{\includegraphics*[height=1in,width=.75in]{../Figures/ParabSecTan}}
%   \label{fig:ParabSecTan}
% \end{wrapfigure}
  Suppose that the two blue line segments in the figure at the right
  are parallel. Show that $c$ is exactly halfway between $a$ and
  $b$. \\
  \comment{That is, show that $c=\frac{a+b}{2}$.}\\
  \centerline{\includegraphics*[height=2in,width=1.5in]{../Figures/ParabSecTan}}  
  % \begin{wrapfigure}[]{r}{.75in}
  %   \vskip-8mm{}
  %     \captionsetup{labelformat=empty}
  %     \centerline{\includegraphics*[height=.5in,width=.75in]{../Figures/ParabSecTan}}
  %     \label{fig:ParabSecTan}
  %   \end{wrapfigure}
    % show that if the tangent line at $C$ is parallel to
    % $\overline{AB}$ then the $x$ coordinate of $C$ is the midpoint of
    % the interval between the $x$ coordinates of $A$ and $B$.
  \end{embeddedproblem}

  \begin{embeddedproblem}[Find the Pattern]
    \label{problem:HyperbolaTriangles}
    Suppose that  $a$ and $n$ are positive numbers. In the following
    sketch the line $AB$ is tangent to the curve $x^ny=1$ at the point
    $\left(a, \frac{1}{a^n}\right)$.\\
      \centerline{\includegraphics*[height=1.4in,width=2.9in]{../Figures/ConstantTriangle12}}
    \begin{enumerate}[label={\bf{}(\alph*)}]{}
    \item For $n=1$ first show that the  coordinates of points  $A$ and $B$ are
        $(0, 2/a)$ and $(2a, 0)$, respectively, and then confirm that:
      \begin{enumerate}[label={\bf{}(\roman*)}]{}
        \begin{multicols}{2}
        \item area$(\triangle AOB)=2$
        \item area$(\triangle OPQ)=1/2$
        \item area$(\triangle BPQ)=1/2$
        \item area$(\triangle PAR)=1/2$ 
        \end{multicols}
      \end{enumerate}
    \item For $n=2$ first show that the  coordinates of points  $A$ and $B$ are
        $\left(0, \frac{3}{a^2}\right)$ and $\left(\frac32a,
          0\right)$, respectively, and then confirm that:
      \begin{enumerate}[label={\bf{}(\roman*)}]{}
        \begin{multicols}{2}
        \item area$(\triangle AOB)=\frac{9}{4a}$.
        \item area$(\triangle OPQ)=\frac{1}{2a}$.
        \item area$(\triangle BPQ)=\frac{1}{4a}$.
        \item area$(\triangle PAR)=\frac{1}{2a}$.
        \end{multicols}
      \end{enumerate}
    \item For $n=3$ first show that the  coordinates of points  $A$ and $B$ are
        $\left(0, \frac{4}{a^3}\right)$ and $\left(\frac43a,
          0\right)$, respectively and then confirm that:
      \begin{enumerate}[label={\bf{}(\roman*)}]{}
        \begin{multicols}{2}
        \item area$(\triangle AOB)-\frac{8}{3a^2}$.
        \item area$(\triangle OPQ)=\frac{1}{2a^2}$.
        \item area$(\triangle BPQ)=\frac{1}{6a^2}$.
        \item area$(\triangle PAR)=\frac{1}{2a^2}$.
        \end{multicols}
      \end{enumerate}
    \item Find a  formula for  the  coordinates of points  $A$ and
      $B$ for any value of $n$ and then:
      \begin{enumerate}[label={\bf{}(\roman*)}]{}
        \begin{multicols}{2}
        \item Find the the area of $\triangle AOB$\\ for any value of $n$.
        \item Find the the area of $\triangle OPQ$\\ for any value of $n$.
        \item Find the the area of $\triangle BPQ$\\ for any value of $n$.
        \item Find the the area of $\triangle PAR$\\ for any value of $n$.
        \end{multicols}
      \end{enumerate}
    \end{enumerate}
  \end{embeddedproblem}

     
\begin{embeddedproblem}
  \label{problem:underl-tang-line}
  Consider the graph of the equation
  $$
  y=x\sqrt[3]{x-8}\\
  $$
  \begin{enumerate}[label={\bf{}(\alph*)}]
  \item Find the equation of the line tangent to the graph of
      when $x=0$.  
    \item Explain what goes wrong when you try to find the equation of
      the tangent line when $x=8$.\\
      \hint{Have you looked at the graph yet?}
    \item Now find the equation of the line tangent to the graph of
      when $x=0$.  
    \end{enumerate}
  \end{embeddedproblem}
  % While it is true that by using Calculus many problems can be solved
  % ``as if by magic'' (to quote Leibniz),
  % Problem~\ref{problem:underl-tang-line} shows that there are also
  % many subtleties involved. We will return to address some of these
  % subtleties later. For now, just remember that things are not always
  % as simple as they might appear.
  


%\pagebreak[4]
These last three problems illustrate that while many otherwise
difficult computations can be done ``in these lines, as if by magic'' (to
quote Leibniz) there are still a great many subtleties to be
accounted for. There is still plenty of room for creative, analytical
problem solving.

We will come back to some of these subtleties later. For now keep in
mind that while Calculus is a very powerful tool it does not allow us
to solve all problems blindly, by simply moving the pebbles.

%\noindent{}\underline{\bf \large The Vomit Comet}\\
%\pagebreak[4]
\section*{The Vomit Comet}
To acclimate astronaut train\-ees to the effects of weightlessness,
NASA uses the following training regimen: It has an airplane perform a
series of steep climbs and sharp dives.  At the top of
\begin{wrapfigure}[10]{r}{2.4in}
%  \vskip-2mm{}
  \captionsetup{labelformat=empty}
  \centerline{\href{https://historycollection.jsc.nasa.gov/JSCHistoryPortal/history/oral_histories/SlezakTR/gallery/pages/S85-44833.htm}{\includegraphics*[height=1.6in,width=2.4in]{../Figures/VomitComet1}}}
  \label{fig:VomitComet}
\end{wrapfigure}
each climb
passengers will experience weightlessness for about $25$ seconds.
During a training flight the pilot will repeat this maneuver about
$40$ times. % NASA can simulate the gravitational force on other bodies,
% such as the moon or Mars as well, by simply adjusting the parameters
% of the flight path.

Because this ``roller coast\-er'' ride sometimes causes nausea for the
passengers the planes used for the maneuver have been christened
``Vomit Comets.''  One of these airplanes was used to film the weightless
scenes in the $1995$ film
\href{https://www.imdb.com/title/tt0112384/}{Apollo 13}.

It is well known that the path of an object moving only under the
influence of the Earth's gravity\aside{That is, we're ignoring air
  resistance.} at the surface of the earth is
parabolic. To simulate weightlessness\aside{More precisely, to create
  neutral buoyancy.} inside the plane the Vomit Comet is flown so that
its flight path matches the parabolic flight path that a thrown object
would follow.

Since the flight path of the Vomit Comet is parabolic we know that its
shape will be the graph of a curve having the form
$$
y=Ax^2+Bx+C,
$$
We'd like to find the values of the unknown parameters $A$, $B$, and
$C$ that match the path of an object moving under the influence of
Earth's gravity.  Once these are known we will also be able to
determine the peak altitude of the flight and how far (horizontally)
it will go before coming back to the altitude where it started the
maneuver. Finally we'd like to confirm the claim that this maneuver takes
about $25$ seconds.

Can we find $A$, $B$, and $C$? No. Unfortunately we don't
yet have enough information to determine $A$, but we can find $B$ and
$C$.

% \pagebreak[4]
\begin{embeddedproblem}
  \label{problem:FlightPath2}
  Suppose the following graph depicts the parabolic flight path
  followed by the Vomit Comet as it starts its maneuver at an altitude
  of $7000$ meters and an initial angle of elevation of $45^\circ$.\\
  \centerline{\includegraphics*[height=2.1in,width=2.1in]{../Figures/VomitComet2}}
  If the equation of this parabola is $ y=Ax^2+Bx+C,$ then  determine the values of $B$ and $C$.
\end{embeddedproblem}

To find $A$ we'll need to have a better model of the motion.
Fortunately, Calculus is exactly the right tool for building such a
model but we don't quite have all of the tools we need yet.  We will
return to the Vomit Comet problem once we have them.


The pilot of a Vomit Comet seeks to simulate weightlessness, but the
pilot of a commercial airliner works hard to \emph{avoid} subjecting
its passengers to extreme effects like weightlessness or, at the other
end, of extreme gravity so it must descend more gradually.  This
situation is a little easier to understand so we will explore it in
the next couple of problems and examples before returning to the Vomit
Comet.

%\pagebreak[4]
\begin{myexample}[Modeling the flight path of an airliner ]
  \label{example:FlightPath}
  We want to model the flight path of a plane as it lands and
  determine the distance from the runway the plane should be when it
  starts its descent. 

  A fundamental tenet of mathematically modeling real world phenomena
  is to keep things as simple as possible. So, the first thing we'd be
  likely to try for is a parabolic descent path:
  $$y=Ax^2+Bx+C.$$
  But it is pretty clear that this won't work because the
  plane should be traveling horizontally at the beginning and at the end
  of its descent. At the end, because at that point it should be on
  the ground, and at the beginning because we don't want to terrify
  the passengers.
  
%\pagebreak[4]
    \begin{ProficiencyDrill}
      Show that there is only one point on any parabola where the line
      tangent to the curve is horizontal. Explain why this proves that the
      flight path cannot be parabolic. \\
      \hint{Wherever the flight path
        is horizontal its slope will be zero.}
    \end{ProficiencyDrill}

  
  The next simplest curve we could use would be a cubic polynomial
  \begin{equation}
    y=Ax^3+Bx^2+Cx+D.\label{eq:FlightPath}
  \end{equation}
  
  \begin{embeddedproblem}
    \label{problem:FlightPath1}
    Below is a section of a cubic polynomial depicting a flight path with the plane
    starting initially at the point $P(l,h)$ and ending at the airport
    $Q$ which we will
    arbitrarily designate as  the origin.\\
    % \centerline{\includegraphics*[height=2.2in,width=4.5in]{../Figures/Descent1}}
    \centerline{\includegraphics*[height=1.5in,width=3.1in]{../Figures/Descent1}}
    \begin{enumerate}[label={\bf{}(\alph*)}]
    \item Assuming that the flight path is the graph of
      Equation~\ref{eq:FlightPath}, compute $\dx{y}$. Then divide by
      $\dx{x}$ to obtain $\dfdx{y}{x}$.
    \item Show that $D=0$. \\
      \hint{The point $(0,0)$ is on the graph
        of equation~\ref{eq:FlightPath}.}
    \item Show that $C=0.$ \\
      \hint{The flight path is horizontal at
        $(0,0)$.}
    \item Determine the values of $A$ and $B$ in terms of $l$ and
      $h$. \\
      \hint{The point $P$ is on the flight path and the flight
        path is horizontal at $P$.}
    \end{enumerate}
  \end{embeddedproblem}
\end{myexample}

%\pagebreak[4]

We've taken our analysis of the flight paths of both the Vomit Comet
and our airliner as far as we can using slopes alone. The problem is
that a flight path is like a road. A road doesn't go anywhere. It just
sits there. But an object traveling along a road (or a flight path) is
moving. It has a velocity and an acceleration.

Analyzing a static path does not allow us to model either the velocity
or the acceleration of an object moving along that path. Fortunately
for us, Newton's formulation of Calculus is just what we need to do
attack these problems.

But since the Vomit Comet aims to replicate the path of a body falling
freely under the influence of gravity  we will also 
need to consider the influence of  gravity. 

%\pagebreak[4]
\section{Galileo Drops the Ball}
\label{sec:rates-change}
\markright{\sc Galileo Drops the Ball}{} 

\begin{wrapfigure}[13]{r}{1.8in}
  \vskip-15mm{}
  \captionsetup{labelformat=empty}
  \centerline{\href{https://mathshistory.st-andrews.ac.uk/Biographies/Galileo/}{\includegraphics*[height=2.in,width=1.8in]{../Figures/Galileo}}}
  \label{fig:GalileoPortrait}
  \caption{\centerline{Galileo Galilei}\\ \centerline{1564-1642}}
\end{wrapfigure}
In the previous section we asserted -- with no justification
whatsoever -- that the flight path of the Vomit Comet would be
parabolic. Then, in Problem~\ref{problem:FlightPath2}, we gave you the
task of finding the constants $B$, and $C$ for which the graph of
the curve $$y=Ax^2+Bx+C$$ is the flight path of the Vomit Comet with
an initial angle of elevation of $45^\circ$. You
should have obtained
$$
y=Ax^2+x+7000.
$$


Inside the airplane, we want what is called \term{neutral buoyancy}.
This means that we don't want any forces pushing us up or down, side
to side, or back and forth \emph{relative} to the
airplane\aside{Naturally, if the plane and everything in it are
  falling towards the earth then the force of gravity is in play. But
  if everything falls together there is no force \emph{inside}
  (relative to) the plane.}. The way to assure this is to force the
airplane to fly along the same path it would take if it were unpowered
and falling freely.  At first it might seem like all we have to do is
shut down the engine, but that won't work because air resistance will
prevent the plane from falling \emph{freely}.  The engines must be
engaged to force the plane along the path it would follow naturally if
there were no atmosphere\aside{Also, it is incredibly dangerous to
  turn off the engines of an airplane while it is flying.}, and thus, no
resistive forces.  In this section we'll see why this path must be
in the shape of a parabola.



Newton very famously said, ``If I have been able to see further, it is
by standing on the shoulders of giants,''
and, in his $1684$ paper on Calculus, Leibniz remarked that ``Other
very learned men have sought in many devious ways what someone versed
in this calculus can accomplish in these lines as by magic.''

As we've seen Fermat, Descartes, and Roberval were three of the
learned giants Newton and Leibniz were indebted to.  Another was
Professor Galilei (1564-1642) of the University of Pisa, who is
universally known and referred to by his first name,
\href{https://mathshistory.st-andrews.ac.uk/Biographies/Galileo/}{Galileo}. As
a professor of mathematics at Pisa, Galileo studied, among other
things, how objects moved under the influence of gravity.

We will be following Galileo's lead to address such questions as:
\begin{itemize}{}{}
\item Suppose you throw a ball straight up into the air. How high will
  the ball go?
\item How long will it take for it to hit the ground and  what will its velocity be at impact?
\item If you throw the ball up  twice as hard. Will it go twice as high?
\end{itemize}

Since Galileo had neither the technology nor the mathematics to
account for air resistance he ignored it.  For the moment we will ignore
it too.

  \begin{embeddedproblem}
    \label{prob:TossedBall1}
    We all know from experience that if you throw a ball straight up
    in the air, it will reach some maximum height. But suppose you
    throw the ball up twice as hard, would it go twice as high?\\
    \comment{We're asking you to guess. Don't worry about being
      wrong. Just take your best guess.}
  \end{embeddedproblem}


  \begin{wrapfigure}[15]{r}{2in}
    \vskip-4mm{}
    \captionsetup{labelformat=empty}
    % \centerline{\href{https://www.youtube.com/watch?v=KDp1tiUsZw8}{\includegraphics*[height=2.6in,width=2.3in]{../Figures/HammerFeather}}}
    \centerline{\href{https://www.youtube.com/watch?v=KDp1tiUsZw8}{\includegraphics*[height=2.3in,width=2in]{../Figures/HammerFeather}}}
    \label{fig:HammerFeather}
  \end{wrapfigure}
  The accepted theory of motion in Galileo's time was
  \href{https://en.wikipedia.org/wiki/Aristotle}{Aristotle}'s
  assertion that a heavier object would fall faster than a lighter
  object. This was an entirely reasonable thing to believe at 
the
  time. Our common experience is that a hammer falls faster than a
  feather. But we now know that this is because the resistance of the
  air slows the feather more than it does the hammer.

  In $1971$, Apollo Astronaut David Scott performed this experiment on
  the moon (see the video at the right). Because there is no air on the
  surface of the moon there is no air resistance. Of course, the
  hammer and feather hit the surface of the moon at the same time.
  Since he didn't have access to the moon Galileo had to be clever
  instead.

  But Galileo was one of the group of new scientists who gathered
  experimental data and applied mathematical principles to
  theories. As a result of his experimental investigations Galileo
  proposed that in the absence of air all objects would fall at the
  same rate\aside{Galileo could be quite argumentative and this
    eventually got him into hot water with the Vatican. In the last
    years of his long life he was threatened with excommunication and
    and risked execution for espousing
    \href{https://mathshistory.st-andrews.ac.uk/Biographies/Copernicus/}{Copernicus'}
    heliocentric theory of planetary motion. It was only because he
    was well known and respected as a scientist that he was only
    confined to his home instead.}.
  
How did he surmise that all objects would fall at the same rate?

\begin{wrapfigure}[10]{r}{2.8in}
  \vskip-3mm{}
  \captionsetup{labelformat=empty}
  \centerline{\includegraphics*[height=1.4in,width=2.8in]{../Figures/GalileoRamp}}
  \label{fig:GalileoRamp}
\end{wrapfigure}
Well actually, he didn't.  We need to be very careful in our use of
language here. To say that ``all objects fall at the same rate'' is a
bit sloppy. Since velocity is the rate of change of position it seems
to say that all objects fall at the same velocity.

But we know that's not true because the velocity of a falling object
depends on how long it has been falling. At the moment  you drop a ball its
velocity is zero. Thereafter it gains
velocity -- it \emph{accelerates.} Thus an object that has just been
dropped is moving slowly while an object that has been falling for a
while is moving at a faster rate.


What Galileo \emph{actually} proposed was that the rate of change of
the velocity -- the \term{acceleration} -- of all falling bodies is
constant.




But an object falling freely under the influence of Earth's gravity
gets moving pretty quickly and in Galileo's day the tools available
for measurement were very limited. So he slowed things down by letting
small balls roll down a ramp assuming (correctly) that friction with
the ramp would not significantly influence his measurements. When he
did this he noticed that the balls always seemed to accelerate at a
constant rate which depended only on the steepness of the ramp. More
precisely, the acceleration of the rolling ball varied with the angle
of descent. In modern terms we would say that the acceleration is a
function of the angle of descent.

\begin{wraptable}[13]{l}{1.9in}
  \vskip-2mm{}
  \captionsetup{labelformat=empty}
  \begin{tabular}{c|c}
    \hline
    $\theta$ \amp{}Acceleration\\
    (in radians)      \amp{} in $(m/s^2)$\\[2mm]
    \hline{}
    $\pi/3\approx1.05$\amp{}4.9\\
    $\pi/4\approx0.79$\amp{}6.93\\
    $\pi/6\approx0.52$\amp{}8.49\\
    $\pi/12\approx0.26$\amp{}9.47\\
    $\pi/24\approx0.13$\amp{}9.72\\
    $\pi/48\approx0.07$\amp{}9.78\\
    $\pi/100\approx0.03$\amp{}9.795\\
    $\pi/200\approx0.02$\amp{}9.799\\
    $\pi/300\approx0.01$\amp{}9.799\\\hline
  \end{tabular}
\end{wraptable}
Referring to the diagram above, it is clear that when\aside{Notice
  that we are using radians to measure the angle. We will do this
  fairly consistently throughout this text because it is usually best
  for scientific computation, But degree measure is common so it will
  occasionally make an appearance, like it did in
  Problem~\#\ref{problem:FlightPath2}. } $\theta=\pi/2$ the ball won't
move at all. So its acceleration is zero. When $\theta=0$ the ball
accelerates downward freely under the influence of gravity with no
resistance\aside{Or at least, very little.}.  Galileo measured the
acceleration associated with steeper and steeper ramps (that is for
$\theta$ closer and closer to zero) obtaining a table of values
similar to the one at the left.

From the table it is clear that as $\theta$ gets closer and closer to
zero the acceleration is getting closer and closer to $9.8.$ Thus
Galileo  deduced that when there is no ramp (when
$\theta =0$) the \emph{velocity} will increase each second by $9.8\,
m/s.$ 

That is, the \emph{velocity} of an object falling under the influence
of the earth's gravity increases by $9.8$ meters per second, each
second. This is usually abbreviated as $9.8\,m/s^2$, is commonly
denoted by $g$, and is called the \term{constant of acceleration due
  to gravity}

\pagebreak[4]
\begin{ProficiencyDrill}
  With a little trigonometry you can deduce that $g=9.8\,m/s^2$
  with only one measurement.  Use the diagram and the first line of the
  table above to deduce that $g=9.8.$\\
  \hint{$\cos\left(\pi/3\right)=1/2$.}
\end{ProficiencyDrill}

All of this means that as an object falls, its \emph{velocity} increases at
the constant rate of \(9.8\, \sfrac{m}{s}\) every second: $v=9.8t.$ If
we drop the object from some height, then its initial velocity is
zero.  After one second it will be falling at a rate of
\(9.8\,\sfrac{m}{s},\) after two seconds, its velocity will be
$9.8\times2=19.6\,\sfrac{m}{s}$, etc.

The specific number that Galileo found, $9.8\,m/s$, is an artifact of
the units we use to measure distance (meters) and time (seconds) and
the fact that we are on the surface of a particular planet.  If we
change our units and measure distance in feet instead then at the
surface of the earth a falling object will accelerate at $32 ft/s^2.$
If we go to  the surface of the moon, it will accelerate at
$1.6\, m/s.$


%\pagebreak[4]
%\begin{samepage}
  \begin{ProficiencyDrill-1line}{}
    Show that if we measure distance in feet, the acceleration constant
    on the moon is approximately $5.2 \frac{\text{feet}}{\text{second}^2}.$
  \end{ProficiencyDrill-1line}
%\end{samepage}
  
The general situation is this: If a falling object's velocity, in meters per second
($m/s$), is changing at a constant rate of $r$ meters per second per
second ($m/s^2$), and $t$ is the number of elapsed seconds then
$ v=rt.$ From Galileo's work we know that at the surface of the Earth
$r=g=9.8$ so
\begin{equation}
  v=9.8t
 \label{eq:Velocity}
\end{equation}

By looking at the units of measurement we can check that this is
reasonable. Since the acceleration, $r,$ is measured in $m/s^2$ and
$t$ is measured in seconds, when they are multiplied a symbolic
cancellation gives
$$
\frac{m}{s^2}\cdot s = \frac{m}{s\cdot \cancel{s}}\cdot \cancel{s} = \frac{m}{s}
$$
or meters per second; the units used to measure
velocity. Notice that Equation~\ref{eq:Velocity} is independent of the weight of the
falling object. Under the influence of gravity alone all objects
accelerate downward at the
same rate.




Galileo's hypothesis that all objects falling solely under the
influence of Earth's gravity accelerate at the same rate came from
both his experimental evidence and a famous thought experiment he
described in his book \pubtitle{On Motion}~(1590).  The experiment
runs as follows: Imagine that two objects, one light and one heavy,
are connected to each other by a string and we drop them from  a great
height, say the top
of the Tower of Pisa. If we assume heavier objects do indeed fall
faster than lighter ones the string will soon pull taut as the lighter
object retards the fall of the heavier object.

So the system  taken as a whole will fall more slowly than the
    heavier object alone.
  But the system as a whole is heavier than either individual object.
 So the system taken as a whole
  will fall faster than the heavier object.
This contradiction leads inexorably to the conclusion that our initial
assumption -- that heavier objects accelerate faster -- must be false.

According to legend Galileo tested his hypothesis by dropping balls of
different weights from the top of the Tower of Pisa, but that is
almost certainly pure legend. He never actually did this.  Probably.

\begin{wrapfigure}[15]{r}{1.8in}
  \vskip-6mm{}
\captionsetup{labelformat=empty}
\centerline{\href{https://www.maa.org/press/periodicals/convergence/mathematical-treasure-galileo-s-two-new-sciences}{\includegraphics*[height=2.4in,width=1.8in]{../Figures/GalileoTwoNewSciences}}}
\label{fig:GalileoTwoNewSciences}
\end{wrapfigure}
Galileo set out his ideas about falling objects in his book
\foreign{Discorsi e Dimostrazioni Matematiche intorno à due nuoue
  Scienze} (Discourses and Mathematical Demonstrations Concerning Two
New Sciences) (1638). This was the last of Galileo's many scientific
works. He wrote it toward the end of his life while he was under house
arrest for heresy and forbidden by the Church to discuss the
heliocentric theory of planetary motion. Rather than remain idle
during this time Galileo revisited his unpublished research from when
he was younger and used that as a foundation for his final opus.  The
two sciences referred to in the title were the science of motion,
which became the foundation of modern physics, and the science of
materials and construction, an important contribution to engineering.

Suppose an object falls $56$ meters from the top of the Tower of Pisa to the
ground. Can we determine how fast it is moving when it strikes the ground?

Recall that the formula $v=gt$  tells us the ball's velocity at any
time $t$. So if $t_g$ is the amount of time it takes for the ball to reach the
ground then  the velocity at the
end would be $v=gt_g$. But how do we find out \emph{exactly} when the ball strikes the ground?

To see the difficulty, suppose we dropped the ball from
a height of $9.8$ meters.  Would it take $1$ second for the ball to
hit the ground?

Clearly not.

To reach the ground after $1$ second the ball would have to
\emph{average} $9.8\,m/s$ during the entirety of that first second.
But at first it is not moving at all (velocity $=0$). After one second
its velocity has increased to $9.8 m/s$. So for the entire duration of
that first second the ball's velocity is \emph{less than} $9.8
m/s$. Thus after $1$ second it has not yet hit the ground because it
was never going fast enough to do so.  Exactly how far it \emph{has} fallen
isn't clear.

Reasoning similarly, at the beginning of the next second the ball is
already falling at $9.8\,m/s$ and thereafter it's velocity increases,
so it falls at least $9.8$ meters during the next second, in addition
to the distance it fell during the first. So, while we cannot yet tell
\emph{exactly} how far it falls during either second, we can say that,
(1) it will not hit the ground during the first  second, and (2)  it
definitely  will hit the ground sometime during 
the next second.


If we had an expression for the ball's \emph{position}, $p$, similar
to equation~(\ref{eq:Velocity}), we would be in much better shape.
Galileo was able to determine the distance the ball fell without the
tools of Calculus, but since learning to use those tools is why we're
here, we will use them.


\section*{\underline{Bringing in Calculus}}

Let $p(t)$ represent the ball's position at time $t.$ That is, $p(t)$
tells us how far the ball has moved from its starting position,
$p(0)$, after $t$ seconds.  Our goal is to find a formula for $p(t)$.

During an (infinitesimal) instant of time $\dx{t}$, the velocity is
virtually constant so during that instant the formula
$$
\text{ position=velocity$\times$time  or }  p=vt 
$$
holds. Differentiating, we see that
$
\dx{p}=v\dx{t}$.
But from equation~(\ref{eq:Velocity}) we also know that $v=9.8t$, so
$ \dx{p}=9.8t\dx{t}$.

\begin{embeddedproblem}
%  \vskip-8mm{}
Thus finding the position of the ball at any
time reduces to finding an expression for $p(t)$ that satisfies the
differential equation
\begin{equation}
\dx{p}=9.8t\dx{t}.
\label{eq:GalileoDiffeq}
\end{equation}


\begin{wrapfigure}[]{r}{1.5in}
%  \vskip-5mm{}
  \captionsetup{labelformat=empty}
  \centerline{\includegraphics*[height=2.4in,width=1.5in]{../Figures/Pisa1}}
  \label{fig:Pisa1}
\end{wrapfigure}
We begin with an educated guess.  From the Power Rule that
    the differential of a quadratic expression like $k t^2$ will be
    the linear term $2k t\dx{t}$. So $p(t)=k t^2$ seems like a
    reasonable guess.
    \begin{enumerate}[label={\bf{}(\alph*)}]
\vskip-5mm{}
    \item Show that $p(t)=kt^2$ satisfies 
      equation~(\ref{eq:GalileoDiffeq}) when $k=4.9$.
    \item Earlier we reasoned that a ball dropped from a height of
      $9.8$ meters would strike the ground  between one and two
      seconds after being dropped. Use $p(t)$ to find out exactly how long it takes.
    \item How long would it take the ball to hit the ground if dropped
      from the top of the Tower of Pisa?
    \item Evaluate $\eval{\dfdx{p}{t}}{p}{56}$. What does this represent
      physically? \\\hint{Remember that $\dx{p}$ is a change in position
        and $\dx{t}$ is a change in time.}
    \end{enumerate}
  \end{embeddedproblem}



  To summarize, we see that if $p = p(t)=kt^2$ denotes the position of
  a dropped ball at time $t$, then $v =\dfdx{p}{t}=gt$ is its velocity
  and $a =\dfdx{v}{t}=\dfdx{\left(\dfdx{p}{t}\right)}{t}=g$ is its
  acceleration\aside{The notation,
    $\dfdx{\left(\dfdx{p}{t}\right)}{t}$, is clearly very
    awkward. We'll define a better notation in the next section.}, as
  measured by Galileo.  Notice that in our diagram we also introduced
  the variable $h$ which represents the height of the ball from the
  ground. The way we choose to define our variables can have a
  considerable impact on the way we understand a given problem. Here,
  for example, both $p$ and $h$ give us the position of the ball but
  if $p(t)=4.9t^2$, then $p=0$ represents the top of the tower and the
  positive direction is down. Conversely, if  $h(t)=56-4.9t^2$, then $h=0$
  represents ground level and the positive direction is up.

  \begin{embeddedproblem-enumerate}
    \begin{enumerate}[label={\bf{}(\alph*)}]
    \item     Consider a ball moving vertically so that its height at time $t$
      seconds is given by\\
      \centerline{$
      h(t) = 56+49t-4.9t^2$ meters}.
      \begin{enumerate}[label={\bf{}(\roman*)}]
      \item Determine the starting position, starting velocity, and
        acceleration of this ball.  (Remember to go through the two step
        process of differentiating and then dividing by $\dx{t}.$ We
        insist.)
      \item Is ``up'' the positive or the negative direction?
      \item How high will the ball go?
      \item When will the ball hit the ground?
      \item What is the impact velocity?
      \item Describe in words the physical situation being modeled by the
        formula $ h(t) = 56+49t-4.9t^2 \text{ meters}$?
      \end{enumerate}
    \item
      In  general the height of an object falling near the surface of
      the earth is given by \\
      \centerline{$
      h(t)=h_0+v_0t-4.9t^2$ meters}.
      \begin{enumerate}[label={\bf{}(\roman*)}]
      \item Show that the initial height is $h_0$, the initial velocity
        is $v_0$, and the acceleration is $=-9.8$.
      \item For simplicity, assume $h_0=0$.  How high will the ball go?
      \item Suppose we double our initial velocity.  Would the ball go
        twice as high?  Compare this to your guess in
        Problem~\#\ref{prob:TossedBall1}?
      \end{enumerate}
      \vskip-6mm{}
    \end{enumerate}
  \end{embeddedproblem-enumerate}
    
Since the graph of $h(t)=h_0+v_0t-4.9t^2$ is a parabola it is tempting to
conclude that we have shown that an object falling under the influence
of gravity alone must be a parabola. But this not correct. We have
consistently assumed that the ball is falling vertically so its flight
path is a straight vertical line, not a parabola. The graph of $h(t)$
is not the flight path of the ball.

Before we proceed we will need to introduce some terminology and
notation. This will help us streamline our discussions.

%\pagebreak[4]
\section{The Derivative}
\label{subsec:Derivative}
\markright{\sc{}The Derivative}
  \begin{wrapfigure}[18]{r}{2in}
    \captionsetup{labelformat=empty}
    \vskip-2mm{}
  \centerline{\href{https://mathshistory.st-andrews.ac.uk/Biographies/Lagrange/}{\includegraphics*[height=2.4in,width=2in]{../Figures/LagrangePortrait}}}
  \label{fig:LagrangePortrait}
  \caption{\centerline{Joseph Louis Lagrange}\\ \centerline{1736-1813}}
\end{wrapfigure}
Differentials are helpful for learning and using the differentiation
rules, but otherwise they are not very useful. Moreover as we have
frequently pointed out, using differentials brings up certain logical
and philosophical questions (Like, ``What are they?'') that are
very difficult to address. When we come
back to these matters in
Chapter~\ref{chapt:whats-wrong-with} you'll see what we mean. For all
of these reasons and more the concept of differentials should be
regarded as a ``convenient fiction'' and nothing more.


On the other hand as we've just seen the \emph{ratio} of
differentials, $\dfdx{y}{x}$, is extremely useful as it can represent
either the slope of a graph or a velocity, depending on the
context. Moreover this ``differential ratio'' is an ordinary real
number so once we get past the underlying idea of differentials, there
are no philosophical problems inherent in its use.

Because of their concerns regarding the validity of using
differentials mathematicians in the $18$th and $19$th centuries, had a
strong motivation to skip over the differential concept entirely and
jump immediately to the (finite) differential ratio in order to make
Calculus more rigorous.
\href{https://mathshistory.st-andrews.ac.uk/Biographies/Lagrange/}{Joseph
  Louis Lagrange} ($1736-1813$) was one of the mathematicians who
tried to do this.  In his $1797$ work
\href{https://www.maa.org/press/periodicals/convergence/mathematical-treasure-lagrange-on-calculus}{\foreign{Th\'{e}orie
    des Fonctions Analytiques}} (\pubtitle{The Theory of Analytic
  Functions}) he even coined a new term and a new notation for
$\dfdx{y}{x}$. He called it the \foreign{fonction d\'{e}riv\'{e}e}
(meaning a function \emph{derived} from another function) and he
replaced the differential quotient $\dfdx{y}{x}$ with the more modern
\term{function} notation $y^\prime(x)$ (pronounced ``$y$ prime of
$x$'').

His attempt to make Calculus rigorous was very clever, but ultimately
unsuccessful so we will not say much about it here. But we will adopt
his terminology and eventually, his notation. The 
\foreign{fonction d\'{e}riv\'{e}e} has come into English as ``derived
function,'' or \term{derivative} for short.  The following are all
\term{derivatives}:
\begin{align*}
  y^\prime(x) \amp{}=\dfdx{y}{x} \text{ (the derivative of $y$ with respect
                to horizontal position $x$)},\\
  z^\prime(t)\amp{}=\dfdx{z}{t} \text{ (the derivative of $z$ with
               respect to time, $t$) }
               \intertext{and in general,}
               \sigma^\prime(\eta)\amp{}= \dfdx{\sigma}{\eta} \text{ (the derivative of $\sigma$ with
                                    respect to $\eta$)}.
\end{align*}

\begin{Digression}[Function Notation, and Prime Notation]
  What we have been denoting as $\dfdx{y}{x}$, Lagrange called a
  ``derived function'' (\foreign{fonction d\'{e}riv\'{e}e}), and he
  denoted it with his prime notation: $y^\prime(x)$.

  OK, but what does that mean?

The ``derived'' part seems clear enough. If
  we have $y=x^2$ then use our differentiation rules to derive
  $y^\prime(x)=\dfdx{y}{x}=2x$. But Lagrange called this a derived
  \emph{function}.

  The function concept is deeply abstract. But it underlies
  Lagrange's ideas, the notation he used and most of modern
  mathematics so it is worth taking a moment to examine the idea
  closely.

  In our experience Calculus students frequently have only a very
  fuzzy notion of the function concept. And this is OK.  It took us
  (mathematicians) many tries over the course of centuries to pin down
  exactly what we mean by ``function.'' It seems only fair that we
  should give our students at least a few weeks to internalize the
  concept.

  Loosely speaking, a function is a rule that connects things. We (the
  authors) know that you don't believe us but it's true. Fundamentally
  that's all a function is. The reason you don't believe us is that
  you probably have many years of frustrating experience with
  functions and this has convinced you that they are very complex and
  abstruse entities. This is also true. Individual functions can be
  extraordinarily complex. But the function \emph{concept} is
  reasonably straightforward: A function is a rule that connects two
  things.

  When Lagrange called $y^\prime(x)=2x$ a \emph{derived function} he
  was saying that $y^\prime$ is the name of a function.  It is
  reasonable then to ask, ``What are the things that $y^\prime$
  connects?''

  But this is easy. $y^\prime$ connects
  \begin{align*}
    1 \amp\text{ with }2, \\
    2 \amp\text{ with } 4,\\
    1.5 \amp\text{ with }2.25,\\
    -3  \amp\text{ with } -6,
  \end{align*}
  and so on. In general, if $x$ is any real number then $y^\prime$
  connects $x$ with two times $x$. Of course we are far too lazy to
  write all of this out explicitly (also it is impossible) so
  we have invented some notation to capture the idea. The symbol
  $y^\prime(x)$, means ``the thing that $y^\prime$ connects $x$ to''
  and the equation
  $$
  y^\prime(x) = 2x
  $$
  says that ``the thing that $y^\prime$ connects $x$ to''
  is\aside{``Is'' means ``equals.''} $2x$.

  So $y^\prime$ truly is a function. An often used metaphor is that of
  the  buttons, or keys\aside{Often called ``function'' keys.}, on a
  calculator. The ``$x^2$'' button implements the function $y(x)=x^2$
  ($y$ connects $x$ with $x^2$.) In this metaphor $x$ is called the
  input and $y(x)$ is called the output.

  But Lagrange called $y^\prime$ a \emph{derived} function. Derived from
  what? Clearly it is derived from the formula $y=x^2$.

  But hold on.  The formula $y=x^2$ says that there is a connection (a
  functional relation) between $x$ and $y(=x^2)$ doesn't it?  If we
  need to choose a name for this function the easiest choice to make
  is $y$.  If we think of $y$ as the function $y(x)=x^2$ then clearly
  \emph{the function} $y^\prime$ is derived from \emph{the function}
  $y$ so we call $y^\prime$ ``the derivative'' of $y$.

  The only thing that \emph{really} changes in all of this is how we
  think about the symbol ``$y$''. Both $y=x^2$ and $y(x)=x^2$ say that
  if $x$ is known we get $y$ by squaring $x$. In the former case
  we think of $y$ as a variable like $x$ and the relation between $x$
  and $y$ is given by the formula $y=x^2$. In the latter we think of
  $y$ as the name of that relation: $y(x)=x^2$.

  Sometimes it is helpful to think one way, and sometimes the
  other. When we  think of $y$ as a variable we will suppress
  the ``$(x)$'' symbolism and write $y=x^2$. When we think of it as a
  function we will write ``$y(x)$.'' In fact, if you look back through
  the earlier chapters of this text you'll see that we've been doing
  that all along.
\end{Digression}

In some contexts Lagrange's prime notation has several advantages over
the differential notation we've been using. Over time it has become
the most common notation for the derivative in mathematics. But the
fact that it took over $100$ years to develop suggests that something
more than just notation is in play here. Since our current task is to
master the differentiation rules we will stick (primarily) to
Leibniz's differential notation as much as possible because using
multiple equivalent notations can be very confusing, especially in the
beginning and we don't want make this harder than it needs to be.


But there will come a time when Lagrange's prime notation will be much
more convenient.  At that point we will casually identify the two
expression $\dfdx{y}{x}$ and $y^\prime(x)$ and we will think of them
both as a function \emph{derived from} the function $y(x)$.

When we do that the differential notation we're currently emphasizing
will take on a slightly schizophrenic aspect. On the one hand
$\dfdx{y}{x}$ represents a ratio of differentials. As such $\dx{y}$
and $\dx{x}$ are distinct (infinitesimal) quantities. On the other
hand $\dfdx{y}{x}$ is the name of a function. As such it is all one
symbol and you cannot detach the pieces any more than you can delete
the letter ``n'' from the name of the sine function; si$(x)$ has no
meaning.

Eventually the differentials we've been using so casually all along
will become a guilty secret. Given $y=y(x)$ we'll use them as a
helpful computational aid (a ``convenient fiction'') while we
\emph{compute}. But as soon as we have
$\dfdx{y}{x}$ in hand we will view it as a single, complete symbol
representing the derivative function, often simply replacing it with
$f^\prime(x)$ as if we are ashamed of having used differentials at all.

As we said, this more advanced viewpoint will come later, but to give
you a preview, consider the following.

Recall that in Descartes' Method of Normals, we had to find a double
root of a polynomial.  To deal with this problem,
\href{https://mathshistory.st-andrews.ac.uk/Biographies/Hudde/}{Johann
  van Waveren Hudde} ($1628$-$1704$) developed an algebraic tool for
determining such double roots. Because this was before the invention
of Calculus it was difficult to see why Hudde's Rule worked. Calculus
allows an alternate derivation of his technique which is much easier
to follow\aside{Actually Hudde did considerably more than this. His
  method for finding double roots was a part of what has been called
  the ``Lost Calculus'' of algebraic functions. If you are interested
  you can read more about this in Jeff Suzuki's award winning article
  \pubtitle{\href{https://www.maa.org/sites/default/files/pdf/upload_library/22/Allendoerfer/suzuki339.pdf}{The
      Lost Calculus (1637-1670): Tangency and Optimization without
      Limits}}}.

  \begin{embeddedproblem}[Hudde's Rule:] Consider any polynomial
    $p(x)=a_0+a_1 x+\cdots+a_n x^n$. Let $a$ and $b$ be any real
    numbers and form the following ``Hudde Polynomial.''
    $$
    {p} ^ {*} \left (x \right ) =a {a} _ {0} + \left (a+b \right ) {a} _
    {1} x+ \left (a+2b \right ) {a} _ {2} {x} ^ {2} + \left (a+3b \right
    ) {a} _ {3} {x} ^ {3} +\ldots+ \left (a+nb \right ) {a} _ {n} {x} ^
    {n}
    $$
    Hudde showed that if $r$ is a double root of $p(x)$, then $r$ is a root of the Hudde
    polynomial $p^*(x).$
    % If you go back and look, the special case in
    % Section~\ref{sec:descartes-normals} was when $a=0$ and $b=1$.
    \begin{enumerate}[label={\bf{}(\alph*)}]{}
    \item Show that if $r$ is a double root of the polynomial $p(x)$
      then it is a root of $p^\prime(x)=\dfdx{p}{x}$.  \hint{If $r$ is a
        double root of $p(x)$, then $p(x) = (x-r)^2q(x)$ for some
        polynomial $q(x)$.}
    \item Show that $p^*(x) = ap(x)+bxp^\prime(x)$ and use this to prove
      Hudde's Rule.
    \end{enumerate}
    \vskip-4mm{}
  \end{embeddedproblem}


  Again, we will (mostly) refrain from using Lagrange's prime notation
  until the differential notation we've been using becomes
  inconvenient. But we will begin to refer to a differential ratio
  like $\dfdx{y}{x}$ as the derivative of $y$ with respect to $x$.
  For now then, a derivative is the result of a two-step process:
  \begin{center}
    \begin{minipage}{.85\linewidth}
      \begin{description}
      \item[First:] \alert{Differentiate to produce differentials.}
      \item[Second:] \alert{Divide through by one differential to
          produce a (finite) derivative.}
      \end{description}
    \end{minipage}
  \end{center}
As you become more proficient you will naturally start to combine
these into a single process. But whenever a differentiation is
particularly complex you will find it helpful to break it down into
simple steps using the differentiation rules you're learning.


It should be clear that the derivative \emph{function} is by far the
more important concept,
% This is one reason we keep referring to
% differentials as ``convenient fictions.'' We only really need
% differentials as stepping stones on our way to building the
% derivative. To be sure, we will continue to use them because as a
% computational aid differentials are indeed very convenient
% fictions.
so more and more we'll need to \emph{think} of the expression
$\dfdx{y}{x}$ as the name of a function. This can be very confusing
because the differential notation was never designed, or intended,
to represent functions.

We will take a moment to address this before we proceed.

If the symbol $\dfdx{y}{x}$ is the name of the function we get by
differentiating $y$ then a natural question is, ``What is the
derivative of the function $\dfdx{y}{x}?$'' In the previous section we
wrote this as $\dfdx{\left(\dfdx{y}{x}\right)}{x}$ but that is clearly
a very clumsy formulation. We'll need something simpler.

The first time we differentiate $y = x^3$ we get its derivative,
$3x^2$ The second time we differentiate $y = x^3$ we get the
derivative of $3x^2$, or $6x.$ What should we call $6x?$ The obvious
choice is to call $6x$ the \emph{second derivative} of $y=x^3.$
Clearly then, $3x^2$ is the \emph{first derivative\aside{But we will
    continue to call it ``the derivative'' when there is no chance
    of confusion.}} and $6$ is the \emph{third derivative} of
$y=x^3.$

All well and good, but it will be cumbersome -- and we are far too
lazy -- to write out ``the third
derivative'' every time we need it. What notation should we use to
indicate higher order derivatives?

For better or worse the extension of the differential notation to
higher order derivatives that has  become standard is this:
If $y = x^3$  then the (first) derivative of $y$ is \\
\centerline{  $\dfdx{y}{x}= 3x^2$}
as before, and the second derivative is\\
\centerline{ $ \dfdxn{y}{x}{2} = 6x.$}
Differentiating one more time we get the third derivative,
$  \dfdxn{y}{x}{3} = 6$.

Notice the inconsistent placement of the exponents.  This is due to
Leibniz. He would have proceeded as follows. Let $y= x^3$. Then\\
\centerline{$\dx{y}= \dx{(x^3)}= 3x^2\dx{x}$.}
So far, so good. Next apply the
Product Rule and we get
\begin{equation}
  \dx{(\dx{y})}= \dx(3x^2\dx{x}) = 3x^2\underbrace{\dx{(\dx{x})}}_{=0} +
  6x\dx{x}\dx{x}\label{eq:SecDiffProblems}
\end{equation}
so
\begin{equation}
  \dx{(\dx{y})} =    6x\dx{x}\dx{x}.    \label{eq:PRNotation} 
\end{equation}

To see why $\dx{(\dx{x})}=0$ we need to know that Leibniz always
considered the differential $\dx{x}$ to be a constant. And the
differential of any constant is zero, by the Constant Rule.


To see why the exponents are placed as they are we note that a
natural way to abbreviate the left side of
equation~(\ref{eq:PRNotation}) is $\dx{(\dx{y})}=\dx^2y$.  On the right
side the expression $\dx{x}\dx{x}$ is the product of $\dx{x}$ with
itself so it is equal to $\dx{x}^2$. Thus when we divide both sides of
equation~(\ref{eq:PRNotation})  by
$\dx{x}^2$ we get $\dfdxn{y}{x}{2}=6x.$  Similarly we get 
$\dfdxn{y}{x}{3} = 6$ and  $\dfdxn{y}{x}{4} = 0$.

The advantage of Lagrange's prime notation should be clear. For higher
order derivatives we simply use more primes. If $y=f(x)$ then
\begin{align*}
  \dfdx{y}{x} \amp{}= f^{\prime}(x),\\
  \dfdxn{y}{x}{2} \amp{}= f^{\prime\prime}(x),\\
  \dfdxn{y}{x}{3} \amp{}= f^{\prime\prime\prime}(x)
\end{align*}
and so on. Above the third derivative Lagrange's primes become
troublesome as well, but that is a discussion for another time.

Obviously we can continue differentiating as many times as we would
like, but this would be rather dull for this particular example since
every derivative after the fourth is zero. The following example is
more interesting.

\begin{myexample}
  \label{example:HighOrderDerivs}{}
  Consider the expression $y=\frac1x =x^{-1}$. As we've seen
  $\dfdx{y}{x}= -1x^{-2}
  =\frac{1}{x^2}.$ Continuing we have
  \begin{align*}
  \dfdxn{y}{x}{2} \amp{}= (-1)(-2)x^{-3},\\
    \dfdxn{y}{x}{3} \amp{}= (-1)(-2)(-3)x^{-4},\\
    \intertext{and finally}
    \dfdxn{y}{x}{4} \amp{}= (-1)(-2)(-3)(-4)x^{-5}
  \end{align*}
  It is an unfortunate fact that you've probably been taught all of
  your life to ``simplify'' complex looking expressions like
  $(-1)(-2)(-3)(-4)$. So much so that you probably do it reflexively,
  without thinking, and are wondering why we left the coefficients
  above in the form we did.  

The reason is simple. We were looking for patterns not numbers.
Writing the above formulas as $\dfdxn{y}{x}{2} = 2x^{-3},$
$\dfdxn{y}{x}{3} = -6x^{-4}$, and $\dfdxn{y}{x}{4} = 24x^{-5}$ would
have obscured the pattern.
Keep this in mind as you proceed.  Algebraic or arithmetical
``simplifications'' often simply get in the way. Don't do them until
there is a compelling reason to do so.

% \pagebreak[4]

\begin{ProficiencyDrill-1line}
    Find the pattern in example~\#\ref{example:HighOrderDerivs} and
    compute $\dfdxn{y}{x}{50}$ directly, without computing all fifty
    derivatives.
  \end{ProficiencyDrill-1line}
\end{myexample}


\begin{myexample}
  Consider the circle $x^2+y^2=1$. Differentiating, we have
  $
  2x\dx{x}+2y\dx{y}=0$, or
  $\dfdx{y}{x}=-\frac{x}{y}$.
  Differentiating again we have
  $$
  \dx{\left(\dfdx{y}{x}\right)}=-\frac{y\dx{x}-x\dx{y}}{y^2}.
  $$
\end{myexample}

  \begin{embeddedproblem}
    Continue this example by showing that
    $  \dfdxn{y}{x}{2}=-\frac{1}{y^3}$.
    Now solve for $y=\pm\sqrt{1-x^2}$ and use this to compute
    $\dfdxn{y}{x}{2}$. Do you get the same answer? Which method do you
    prefer?
  \end{embeddedproblem}

%\pagebreak[4]
\begin{ProficiencyDrill}
  For each of the following find $\dfdxn{y}{x}{2}$ in terms of $x$
  and $y$:
  \begin{enumerate}[label={\bf{}(\alph*)}]
    \begin{multicols}{2}
    \item $y=3x^4-x^3+2x-7$
    \item $x=y^2$
    \item $y=\sqrt{x}$ Compare this with part (b).
    \item $xy=1$ Compare this with
      Example~\#\ref{example:HighOrderDerivs}. Which method do you
      prefer?
    \item $\frac{x^2+y}{3x+y^2}=x-y$ 
   \end{multicols}
  \end{enumerate}
\end{ProficiencyDrill}


  We know that it is not generally true that $a^b=a\cdot b$ even
  though there are certain exceptions, like $a=b=1$, $a=4$ and
 $b=1/2$, or $a=b=2$. In the same way, even though the Product Rule
  makes it very clear that
  \begin{equation}
    \dfdx{(y\cdot z)}{x}\neq  \dfdx{y}{x}\cdot\dfdx{z}{x}.\label{eq:NewbyError1}
  \end{equation}
  there are certain pairs of functions which are exceptions; for which
  equation~(\ref{eq:NewbyError1}) \emph{is} true.

\pagebreak[4]
\begin{embeddedproblem}[Find the pattern]
  \label{problem:NewbyError1}
  For example, show that for each of the following it \emph{is} true
  that $
  \dfdx{(y\cdot z)}{x}=  \dfdx{y}{x}\cdot\dfdx{z}{x}$.
  \begin{enumerate}[label={\bf{}(\alph*)}]
  \item 
    \begin{enumerate}[label={\bf{}(\roman*)}]
      \begin{multicols}{3}
      \item
        $y=x$\\
        $z=\frac{1}{1-x}$
      \item
        $y=x^2$\\
        $z=\frac{1}{(2-x)^{2}}$
      \item
        $y=x^3$\\
        $z=(3-x)^{-3}$
      \end{multicols}
    \end{enumerate}
  \item Find the general pattern in part (a).
  \item Those pairs of functions which fit the pattern you found in
    part (b) are not the only exceptional pairs. Can you find others?
  \end{enumerate}
\end{embeddedproblem}


\begin{mynotation}{Some Troubling Aspects of Differentials and Differential}
  % \begin{mynotation}{Function}
  We would be remiss if we did not point out  some very troubling
  aspects of the differential notation, especially when we extend it
  to  higher order derivatives. 

  One very problematic and glaring question is, ``If $\dx(\dx{x}) = 0$
  why isn't $\dx(\dx{y})$ also equal to zero?'' That is a much deeper
  question than it appears to be. We will address this in
  Section~\ref{sec:curvature}.

  But there are other problems as well.  When we introduced the
  concept of a differential in Chapter \ref{chapt:differentials} we
  appealed to your understanding of \emph{finite} differences.
  That is, if $x_1\lt{}x_2$ then the finite difference
  $$
  \Delta{}x=x_2-x_1
  $$
  is a measure of the actual separation between $x_2$ and $x_1.$
  Extending this to infinitesimals,  we imagined that $x_2$ and $x_1$ are
  \emph{infinitely} close together -- but still distinct -- and we
  represented the distance between them as
  $$
  \dx x = x_2-x_1.
  $$
  % It is all well and good to \emph{write down} expressions like $\dx y$
  % or $\dx^2 y$ but what do they represent in reality?
  
  If $x_1$ and $x_2$ are infinitely close together doesn't that mean
  that they are not separated at all? If they are not separated then
  $x_2=x_1.$ In that case isn't $\dx x =x_2-x_1$ just zero? And if
  $\dx x =0$ then doesn't the expression $\dfdx{y}{x}$ call for
  division by zero? Which we know is meaningless?

  This is all very troubling, which is why we have encouraged you to
  not think about it very much so far. But as higher order derivatives
  come into play things begin to get truly weird and it is time to
  \emph{begin} thinking about it. A little.

  Consider the actual meaning the second derivative must have in terms
  of differentials. If the first derivative is obtained by computing
  the ratio of quantities that are infinitely close together, what can
  it possibly mean to differentiate it again? If we are using numbers
  that are infinitely close together to compute the first derivative,
  what can we possibly be using to compute the second, or the third
  derivative?

  What is the infinitesimal difference of the infinitesimal difference
  of an infinitesimal difference?

  Reacting to criticism of religious
  mysteries from the scientific community one early critic\aside{This
    is Bishop George Berkeley, from whom we will hear a great deal
    more in chapter~\ref{chapt:whats-wrong-with}. } of Calculus put it
  like this,
\begin{center}
  \begin{minipage}{.75\linewidth}
    \it{}
    ``But he who can digest .\ .\ .\  a second or third
  Difference, need not, methinks, be squeamish about any Point in
  Divinity.''
\end{minipage}
\end{center}

This is all very puzzling and it is tempting to toss it all away.
Except that at this point we have a great deal of evidence that if
$y=x^3$ then the formula
$$
\dfdxn{y}{x}{2}=6x
$$
is both quite meaningful and quite useful. There must be something to
this. But what?

  If we are being honest with ourselves we must acknowledge the
  legitimacy of these concerns. There is something here that we do not
  yet understand. We will proceed very cautiously.

  In particular we will avoid entirely the very delicate issue of
  trying to take a the differential of a differential. That is, we
  will only differentiate finite expressions. For example if $y=x^3$,
  then $\dx{y}=3x^2\dx{x}$ so $\dfdx{y}{x}=3x^2$ which is clearly
  finite. Differentiating again gives
  $\dx{\left(\dfdx{y}{x}\right)}=6x\dx{x}$, so
  $$
  \dfdxn{y}{x}{2}=\dfdx{\left(\dfdx{y}{x}\right)}{x}=6x.
$$
Any attempt to differentiate a differential cuts very close to the
foundational issues we've alluded to, and have been very careful to
avoid so far. It is not yet time for us to consider these questions so
we will leave them for now. But that time is coming.
\end{mynotation}

Adopting Lagrange's terminology, but not his notation, we see that if
the position of a point moving in a straight line (like the $x$ axis)
is given by $x=x(t)$, then the first derivative, $\dfdx{x}{t}$, will
give its velocity, and its second derivative, $\dfdxn{x}{t}{2}$, will
give its acceleration.

\begin{ProficiencyDrill}
  Each of the following represents the position of a point on
  the $x$-axis at time $t$.  Find the velocity and acceleration.
  \begin{enumerate}[label={\bf{}(\alph*)}]
    \begin{multicols}{2}
    \item $x(t)=12t^3$
    \item $x(t)=-4t^4+3t^2+1$
    \item $ x(t)=5-2\sqrt{t}+t^3$
    \item $\displaystyle x(t)=\frac12t^{1/2}+t^{-1/2}$
    \item $\displaystyle x(t)=\frac{1}{\sqrt{t^2+t+1}}$
    \item $x(t)= t^{2/3}$
    \end{multicols}
  \end{enumerate}
\end{ProficiencyDrill}

\pagebreak[4]
\begin{embeddedproblem}
  For each of the following, determine if a point moving along the
  $x$-axis is slowing down or speeding up at the instant
  $t_0$.
  \begin{enumerate}[label={\bf{}(\alph*)}]
    \begin{multicols}{2}
    \item $\left.\dfdx{x}{t}\right|_{t=t_0}\gt{}0, \left.\dfdxn{x}{t}{2}\right|_{t=t_0}\gt{}0$ 
    \item $\left.\dfdx{x}{t}\right|_{t=t_0}\gt{}0, \left.\dfdxn{x}{t}{2}\right|_{t=t_0}\lt{}0$ 
    \item $\left.\dfdx{x}{t}\right|_{t=t_0}\lt{}0, \left.\dfdxn{x}{t}{2}\right|_{t=t_0}\gt{}0$ 
    \item $\left.\dfdx{x}{t}\right|_{t=t_0}\lt{}0,
      \left.\dfdxn{x}{t}{2}\right|_{t=t_0}\lt{}0$ 
    \end{multicols}
  \end{enumerate}
%  \vskip-6mm{}
\end{embeddedproblem}

\section{Thinking Dynamically}
\label{sec:newt-flux-dynam}
\markright{\sc{}The Dynamic Approach}

\begin{wrapfigure}[12]{r}{1.5in}
  \vskip-5mm{}
  \captionsetup{labelformat=empty}
  \centerline{\includegraphics*[height=2in,width=1.5in]{../Figures/Helicopter1}}
  \label{fig:Helicopter1}
\end{wrapfigure}
If $p(t)$ represents the position of an object moving in a straight
line we have seen that its velocity is given by $\dfdx{p}{t}$. When
the object is falling vertically we will use $y(t)$ to represent its
vertical position, so that $\dfdx{y}{t}$ is its vertical
velocity. Similarly if it is moving horizontally we will use $x(t)$
and $\dfdx{x}{t}$.


% In we imagined a point moving
%  with a constant horizontal  velocity of $1$ and a  vertical velocity equal
%  to $-9.8$. If the velocities are measured in meters per second this
%  actually models the motion of

 Suppose an  object is 
released  from,  a helicopter  flying horizontally with a
speed of $\dfdx{x}{t}=1\sfrac{\text{m}}{\text{s}}$.  From Galileo's
work we now know that its vertical velocity will be
$\dfdx{y}{t}=9.8 t\, \sfrac{\text{meters}}{\text{second}}$. You may
recall that we posed this more abstractly in
Problem~\ref{problem:HelicopterFall}.

After its release the ball will fall neither horizontally, nor
vertically. It won't even fall along a straight line. The combination
of its horizontal and vertical velocities will cause the flight path
to curve into an arc be similar to that shown in the
figure at the right. We've already stated that this curve will be a
parabola. We are finally in a position to show that this is true.

As the ball falls there are three distinct quantities we are
interested in: the horizontal velocity of the ball, the vertical
velocity of the ball, and its speed in the direction of travel.

The first two we already understand from our work in the previous
section. The horizontal velocity will be $\dfdx{x}{t}$, and the
vertical velocity will be $\dfdx{y}{t}$.

\begin{ProficiencyDrill}
  Since $\dfdx{x}{t}=1$ we have $x=t$. Since  $\dfdx{y}{t}=-9.8t$ we know from our work
  in the previous section that $y=-4.9t^2.$ Use this to find a formula
  for $y=y(x)$ and confirm that the graph is a parabola.
\end{ProficiencyDrill}
% \begin{wrapfigure}[]{l}{.67in}
%   \captionsetup{labelformat=empty}
%   \centerline{\includegraphics*[height=1in,width=.67in]{../Figures/Helicopter2}}
%   \label{fig:Helicopter2}
% \end{wrapfigure}
We'd like to find the \emph{speed} of the object in the direction it
is moving.  Referring to the differential triangle in our sketch it is
clear that this will be\aside{The absolute value bars here are correct
  but don't fret about it too much for now. The reason for this has
  more to do with our vocabulary (speed vs. velocity) than with the
  physical problem. We will discuss this in more depth in
  Digression~\#\ref{digress:SpeedVelocityAbsVal}{} in
  chapter~\ref{cha:calc-trig}} $\abs{\dfdx{s}{t}}$.  But how do we
find $\dfdx{s}{t}$?

Since $\dx{x}$ and $\dx{y}$ are perpendicular (they are displacements
in the horizontal and vertical directions, respectively)  by
the Pythagorean Theorem  we have
$$
(\dx{s})^2=(\dx{x})^2+(\dx{y})^2.
$$
Thus the displacement, $s$, in the direction of travel is
$$
\abs{\dx{s}}=\sqrt{(\dx{x})^2+(\dx{y})^2}.
$$
Dividing by $\abs{\dx{t}}$ we see that the speed in the direction of travel
at any time $t$ is
$$
\abs{\dfdx{s}{t}}=\frac{\sqrt{(\dx{x})^2+(\dx{y})^2}}{\dx{t}} =
\sqrt{\left(\dfdx{x}{t}\right)^2+\left(\dfdx{y}{t}\right)^2}
= \sqrt{1+(9.8t)^2}.
$$

\begin{wrapfigure}[]{r}{2.4in}
\captionsetup{labelformat=empty}
\centerline{\includegraphics*[height=1.2in,width=2.4in]{../Figures/Descent1}}
\label{fig:Descent2}
\end{wrapfigure}
We now hava all of the tools we need to analyze the flight path of our
commercial airliner, seen at the right.

Recall that in Problem~\ref{problem:FlightPath1} we had you find the
equation of this flight path.  The answer turns out to be
$$
y=\frac{-2H}{L^3}x^3+\frac{3H}{L^2}x^2.
$$
But the picture above is static. It is a picture of the path the
airliner has flown after the flight is finished. This tells us very
little about the characteristics of the flight as it occurs.

  \begin{embeddedproblem-enumerate}
    \label{problem:AirlinerSpeed}
    \begin{enumerate}[label={\bf{}(\alph*)}]
    \item Assume that horizontal (ground) speed is maintained at a
      constant $v$ meters/second.  What is the airspeed (speed in the
      direction of travel) of the plane?  In particular, is the
      airspeed also constant?
    \item Show that in general, speed and slope are related by the
      formula:\\
      \centerline{$\displaystyle
        \abs{\dfdx{s}{t}}=\abs{\dfdx{x}{t}}\sqrt{1+\left(\dfdx{y}{x}\right)^2}$.}
    \item Assume that the ground speed is constant and use the
      formula in part (b) to determine where along the flight path the
      plane is traveling the fastest and the slowest.\\      
      \hint{Look at the graph of the flight
        path.}
    \end{enumerate}
  \end{embeddedproblem-enumerate}

\pagebreak[4]
\begin{ProficiencyDrill}
  Suppose  $(x,y)$ are  the coordinates of a ball moving
  along each of the following curves.
    \begin{enumerate}[label={\bf{}(\Roman*)}]
    \begin{multicols}{4}
    \item $y=x^2,$\\ at $(-1,1)$\\ and $(1,1)$
    \item $y=x^3,$\\ at $(-1,-1)$\\ and $(1,1)$
    \item $x^2-y^2=3^2$\\ at $(-5,4)$\\ and $(-5,-4)$
    \item $x^2+y^2=5^2$\\ at $(3,4)$\\ and $(3,-4)$
    \end{multicols}
  \end{enumerate}
If the ball is moving with a constant
  horizontal velocity of
  $\dfdx{x}{t}=2\sfrac{\text{ units}}{\text{second}}$ find each of
  the following.
  \begin{enumerate}[label={\bf{}(\alph*)}]
  \item The vertical velocity of the ball at the points indicated.
  \item The horizontal and vertical acceleration of the ball at  the points
    indicated.
  \item The velocity in the direction of travel of the ball at
    the point indicated.
  \end{enumerate}
\end{ProficiencyDrill}

 The next problem will complete our analysis of the flight path
 of our commercial airline.

%\pagebreak[4]
\begin{embeddedproblem}
  The flight path you obtained in Problem~\#\ref{problem:FlightPath1}
  was $ y=\frac{-2H}{L^3}x^3+\frac{3H}{L^2}x^2 $.
  We need to put some limitations on the vertical acceleration,
  $\dfdxn{y}{t}{2}$, experienced by the passengers in a commercial
  airliner. For simplicity we'll assume that the pilot must maintain a
  constant horizontal speed of $v\sfrac{\text{
      meters}}{\text{second}}$. That is we set  $\dfdx{x}{t}=-v$.)
  \begin{enumerate}[label={\bf{}(\alph*)}]
  \item Use the above equation to show that the vertical acceleration
    is given by
    $$ \dfdxn{y}{t}{2}=\frac{6Hv^2}{L^2}\left(1-\frac{2x}{L}\right).$$
  \item On the interval $0\le x\le L$, what is the largest vertical
    acceleration and what is the smallest vertical acceleration and
    where do they occur?  Does this make sense physically?
  \item Suppose we put a restriction on the vertical acceleration,
    namely
    $$
    -k\sfrac{\text{ meters}}{\text{second}^2}\le\dfdxn{y}{t}{2}\le k\sfrac{\text{ meters}}{\text{second}^2}
    $$
    for some constant $k$.  Show that with this restriction,
    $
    L\ge\sqrt{\frac{6Hv^2}{k}}
    $.
  \item Suppose that initially $H=10000$ meters, $v=100\sfrac{\text{
        meters}}{\text{second}}$,
    $k=.1\sfrac{\text{ meters}}{\text{sec}^2}$ (which is approximately $1\%$ of
    the acceleration due to gravity).  Find what $L$ must be (in
    kilometers).% \\
    % \comment{This represents the minimum distance from
    %   the airport where the plane must begin its descent.}
  \end{enumerate}
\end{embeddedproblem}

 \pagebreak[4]


Earlier we mentioned that NASA claims that the Vomit Comet can
make passengers experience weightlessness for about $25$ seconds.
Let's check on that claim.


To simulate weightlessness (neutral buoyancy) the pilot must
execute a parabolic flight path $y=Ax^2+Bx+C$. In
Problem~\#\ref{problem:FlightPath2} you should have found that $B$
and $C$ were $1$ and $7000$, respectively, so the flight path is
$$
y=Ax^2+x+7000
$$
with $A$ yet to be determined.


The pilot will climb at an angle of $45^\circ$ to an altitude of about
$7000$ meters and then follow this parabolic path to produce a
vertical acceleration of $\dfdxn{y}{t}{2}=-9.8\,\sfrac{m}{s^2}$
(matching the acceleration due to gravity) and horizontal acceleration
of $\dfdxn{x}{t}{2}=0$. This will provide neutral buoyancy inside the
plane.  On the way back down the pilot pulls out of this dive when the
altitude returns to $7000$ meters.

For training purposes this is repeated $40$ times.

%\pagebreak[4]
\begin{embeddedproblem-enumerate}
  \begin{enumerate}[label={\bf{}(\alph*)}]
    % \item What are $y(0)$ and $\eval{\dfdx{y}{x}}{x}{0}$? Use these
    %   to compute $B$ and $C$.
  \item To determine $A$ we need one more fact.  At the beginning of
    the maneuver, the initial airspeed is about $180\sfrac{\text{ meters}}{\text{second}}$
    (approximately $400$ mph).  Use this to determine $\dfdx{x}{t}$
    and in turn use this and the fact that $\dfdxn{y}{t}{2}=-9.8$ to
    determine $A$.
  \item Now that we've determined the flight path, use this to
    determine the value of $x$ is when the pilot pulls out of the dive
    at $7000$ meters. 
  \item Next find the value of is $t$ when the pilot pulls out of the
    dive. How does this compare with the $25$ second claim.
  \item What would the airspeed of the plane be in terms of $t$ ($t=0$
    representing the start of the maneuver) if the pilot wants to
    maintain a constant horizontal speed throughout the maneuver?
  \end{enumerate}
\end{embeddedproblem-enumerate}


\section{Newton's Method of Fluxions}
\label{sec:newt-meth-flux}
\begin{wrapfigure}[10]{r}{1.5in}
  \vskip-20mm{}
\captionsetup{labelformat=empty}
\centerline{\href{https://mathshistory.st-andrews.ac.uk/Biographies/Kepler/}{\includegraphics*[height=2in,width=1.5in]{../Figures/Kepler}}}
\caption{\centerline{Johannes Kepler}\\\centerline{1571-1630}}
\label{fig:Kepler}
\end{wrapfigure}
Through experimentation Galileo had accurately described the motion of
falling objects near the surface of the earth. At about the same time
\href{https://mathshistory.st-andrews.ac.uk/Biographies/Kepler/}{Johannes
  Kepler} ($1571-1630$) had accurately described the motion of the
planets by analyzing the vast catalog of astronomical observation made
by
\href{https://mathshistory.st-andrews.ac.uk/Biographies/Brahe/}{Tycho
  Brahe} ($1546-1601$).  But these two descriptions of motion did not
appear to be related.

It was left to Newton to unify them with a single theory but the
mathematics for this did not exist at the time. So he invented it. He
called his invention
\pubtitle{\href{https://www.maa.org/press/periodicals/convergence/mathematical-treasure-newtons-method-of-fluxions}{The
    Method of Fluxions}}.  
Galileo studied motion on the surface of the Earth.  Kepler studied
the motion of the planets, but the underlying theme for both was
motion.

For Newton everything was in motion.  When he used the variable $x$ he
thought of it as representing something ``flowing in time''
(moving). Such quantities he called \term{fluents}, from
\foreign{fluere}, the Latin word which means ``to flow.''  As he put
it himself in his book
\href{https://www.maa.org/press/periodicals/convergence/mathematical-treasure-newton-on-quadrature-of-curves}{\foreign{Quadratura
    Curvarum}} (On the Quadrature of Curves),
\begin{center}
  \begin{minipage}{.75\linewidth}\it
    I sought a method of determining quantities from the velocities
    of the motions or [of the] increments, with which they are
    generated; and calling these velocities of the motions, or [of
    the] increments, \emph{fluxions}, and generated quantities
    \emph{fluents}, I fell by degrees, in the years $1665$ and
    $1666$, upon the method of fluxions, which I have made use of
    here in the quadrature of curves.
  \end{minipage}
\end{center}


\begin{wrapfigure}[16]{l}{1.5in}
  \vskip-5mm{}
\captionsetup{labelformat=empty}
\centerline{\href{https://mathshistory.st-andrews.ac.uk/Biographies/Brahe/}{\includegraphics*[height=2in,width=1.5in]{../Figures/Brahe}}}
\caption{\centerline{Tycho Brahe}\\\centerline{1546-1601}}
\label{fig:Brahe}
\end{wrapfigure}
For a given fluent $x$, Newton used $\dot{x}$ to refer to its
\emph{instantaneous velocity} or \term{fluxion}. Whereas for Leibniz
the static differential was the fundamental concept, for Newton the
dynamic fluxion (velocity) was fundamental.  On the surface Newton’s
fluxions seem quite different from Leibniz's differentials.

In Newton's view the only independent variable is time. So \emph{all}
fluxions were rates of change with respect to time, or velocities. To
see how they are connected to the differential ratios we've been using
until now we need only ask ourselves how Leibniz would express the
fluxion, $\dot{x}$, of the fluent, $x$. In Leibniz' notation the rate
of change of $x$ \emph{with respect to} $t$ is the differential ratio
$\dfdx{x}{t}$.  Since Newton \emph{defines} $\dot{x}$ to be the rate
of change of $x$ with respect to (t)ime (velocity), clearly
$$
\dot{x} = \dfdx{x}{t}.
$$


From Newton's point of view differentiating the spatial coordinate $y$
with respect to the spatial coordinate $x$ to get $\dfdx{y}{x}$ is
simply not meaningful, and his dot notation reflects this.  On the
other hand, if $P=(x,y)$ is a point moving along some curve in the
plane then both $x$ and $y$ are fluents with corresponding fluxions
$\dot{y}$ and $\dot{x}$. What Leibniz's notation expresses as
$\dfdx{y}{x}$ Newton's notation expresses as $\frac{\dot{y}}{\dot{x}}$
since
$$
\frac{\dot{y}}{\dot{x}} = \frac{\dfdx{y}{t}}{\dfdx{x}{t}}=
\frac{\dx{y}}{\cancel{\dx{t}}}\frac{\cancel{\dx{t}}}{\dx{x}}=\dfdx{y}{x},
$$
but they are the same thing.
                            
Although the two formulations are equivalent Leibniz' notation has
become dominant in mathematics, but  in fields where velocity is a
fundamental concept, like physics and engineering, Newton's dot
notation is often still used. 

For example, suppose we were considering the curve $y=x^2$ and wanted
the slope of the curve at a point $P$ on the curve.  Using our
differentiation rules, we have
$$
\dx{y}=2x\dx{x}.
$$
To get the slope we divide by $d{x}$ to get
$$
\dfdx{y}{x}=2x.
$$

Using his Method of Fluxions, Newton would have considered $P=(x,y)$
to be a point moving along a curve. This makes both $x$ and $y$
fluents.  His version of  Calculus would have started with the same
governing equation $y=x^2$ and determined that their fluxions were
related by the equation
$$
\dot{y}=2x\dot{x}.
$$

To get the slope of the curve $y=x^2$, we simply  need to recognize
that this is the same as
$$\frac{\dot{y}}{\dot{x}}=2x=\frac{\dfdx{y}{t}}{\dfdx{x}{t}}=\dfdx{y}{x}.$$

% We've been using  Leibniz's notation almost exclusively so far and we
% will continue to do so for a while longer. However when a problem is
% dynamical in nature Newton's viewpoint 
% However, when we are looking at objects in motion
% Newton's more dynamical viewpoint is often more natural. We will
% continue using Leibniz' notation going forward because it is used
% exclusively in mathematics.
As we've mentioned Newton's dot notation has fallen out of favor in
mathematics but it is still used in physics and engineering, so you
will likely see it being used in your physics or engineering courses
(if you take any). We would be remiss if we failed to recognize this
so when it is appropriate we will sometimes couch our problems in the 
dynamical language and dot notation of Newton. We will even
sometimes refer to an expression like $\dfdx{y}{t}$ (where $t$
represents time) as a fluxion, because that's what it is. It's just
given in differential notation.

If you prefer the Leibnizian terms and notation it is easy to
translate between Newton and Leibniz. If $x$ is changing in time then
$\dot{x}=\dfdx{x}{t}$ is the fluxion of $x$.


\begin{ProficiencyDrill}{}
  For each of the following equations, find an equation relating their
  differentials and use this to relate their fluxions (instantaneous
  rates of change with respect to time).
  \begin{enumerate}[label={\bf{}(\alph*)}]{}
    \begin{multicols}{2}
    \item $\displaystyle xy^2-7x+\frac{y}{x}=1$
    \item $\displaystyle x^{\frac12}+y^{\frac12}=1$
    \item $\displaystyle \frac{x^2+y}{\sqrt{y}+3x}=x+7$
    \item $\displaystyle \frac{y}{x}=z^2$
    \item $\displaystyle \sqrt{z^3}=x^2+y^2$
    \item $\displaystyle xyz=1$
    \end{multicols}
  \end{enumerate}
\end{ProficiencyDrill}

\begin{embeddedproblem-enumerate}
  \begin{enumerate}[label={\bf{}(\alph*)}]
  \item Suppose the length of a rectangle is increasing at a rate of
    $1 \sfrac{\text{ unit}}{\text{second}}$ and the width is decreasing
      at a rate of $1$ unit/second.  Make a guess: Will the area
      remain constant?
  \item Let the length be denoted by $L$ and the width denoted by
    $W$. The fluxions of $L$ and $W$ are $\dfdx{L}{t}=1$ and
    $\dfdx{W}{t}=-1$, respectively.  If $A$ denotes the area of the
    rectangle, then compute the fluxion $\dfdx{A}{t}$.  How does
    this compare to your guess in part (a)?
  \item Can you translate your solution into Newton's dot notation?
  \end{enumerate}
\end{embeddedproblem-enumerate}

\pagebreak[4]

\begin{embeddedproblem}
  \vskip-6mm{}
  \begin{wrapfigure}[]{r}{1.5in}
    \vskip-6mm{}
\captionsetup{labelformat=empty}
    \centerline{\includegraphics*[height=1.5in,width=1.5in]{../Figures/SlidingLadder}}
\label{fig:SlidingLadder}
\end{wrapfigure}
The sketch at the right represents a $15$ foot long ladder is leaning
against a vertical wall.  Suppose that the bottom is sliding to the
right at a constant rate.
  \begin{enumerate}[label={\bf{}(\alph*)}]
  \item Make a guess: Is the top sliding down at a constant rate?
  \item  Find $\dfdx{a}{t}$ in terms of $\dfdx{b}{t}$, and use this to
    check your guess in part (a).
  \item Can you translate your solution into Newton's dot notation?
  \end{enumerate}
  \vskip-6mm{}
\end{embeddedproblem}

\vskip-5mm{}

\begin{embeddedproblem}{}
  Suppose that the price of a certain commodity is increasing at a
  rate of 5\% and the quantity sold is decreasing at a rate of
  3\%. Would the revenue increase at a rate of 2\%?  Explain.
\end{embeddedproblem}

\vskip-5mm{}

  \begin{embeddedproblem}{}
    The strength of a signal from a cell tower, measured in
    decibels\aside{The icon on your cell phone probably uses bars,
      but there is no universally agreed upon standard for what ``one
      bar'' means. There is a standard for decibels. } is inversely
    proportional to the square of the distance between the tower and
    the cell phone.  Suppose you have a $200$ foot tall cell tower and
    a car driving away from the tower at a rate of $100 \sfrac{\text{
        ft}}{\text{second}}$
    (approximately $68$ mph).
    \begin{enumerate}[label={\bf{}(\alph*)}]{}
      \item Find a formula for the rate at which the signal is
        decreasing in terms of the distance $x$ from the car to the base
        of the tower.  (Use $k$ to denote the constant of
        proportionality.) 
      \item Plot a graph of the formula in part (a) with $x \ge 0$ and
        use this to approximate the distance from the base of the
        tower where the signal decreases fastest. (Use $k = 1000000$.
        You’ll see why we chose this scale once you’ve plotted it.)
      \item Based on the graph, what is happening to the rate at which
        the signal is decreasing as the car moves farther away from
        the tower?
      \end{enumerate}
        \vskip-5mm{}
  \end{embeddedproblem}

  \vskip-3mm{}
  \begin{embeddedproblem}{}
    It is reasonable to assume that the rate at which the volume of a
    spherical raindrop is increasing as it forms is proportional to
    the surface area of the raindrop.  Show that under this
    assumption, the radius (surprisingly?) grows at a constant rate.
  \end{embeddedproblem}

  \pagebreak[4]
  
    \begin{myexample}
  \begin{wrapfigure}[]{r}{4in}
    \vskip-15mm{}
\captionsetup{labelformat=empty}
\centerline{\includegraphics*[height=1.4in,width=4in]{../Figures/PulleyProblem}}
\label{fig:PulleyProblem}
\end{wrapfigure}
\vskip-5mm{}
Consider two carts joined by a $100$ foot length of rope passing
through a pulley $P$, which is held at a fixed height of $10$ feet
above where the rope attaches to the carts.  Suppose cart $A$ is being
pulled toward the left at a constant speed of
$1 \sfrac{\text{ foot}}{\text{second}}.$


  \begin{ProficiencyDrill-1line}
    Make a guess: Is cart $B$ moving at a constant speed?
  \end{ProficiencyDrill-1line}

  \begin{wrapfigure}[8]{l}{3.5in}
    \vskip-6mm{}
\captionsetup{labelformat=empty}
  \centerline{\includegraphics*[height=1.3in,width=3.5in]{../Figures/TetheredCarts}}
\label{fig:TetheredCarts}
\end{wrapfigure}
Let's check your guess.  The general approach is to find the
relationships between all of the fluents involved and differentiate
and divide by $\dx{t}$ to see how their fluxions are related.  At the
left we have redrawn the essential features of the problem
schematically and labeled all of the fluents.

  We have the following relationships between the fluents:
  $$
  x^2+10^2=z^2,\hskip3 mm y^2+10^2=w^2,\text{ and } z+w=100.
  $$
  \begin{embeddedproblem-enumerate}
    \label{problem:TwoCarts}
    \begin{enumerate}[label={\bf{}(\alph*)}]
    \item Differentiate the above equations to obtain relationships
      between the differentials $\dx{x}$, $\dx{y}$, $\dx{z}$, and
      $\dx{w}$.
    \item Determine the relationship between the fluxions
      $\dfdx{x}{t}$ and $\dfdx{y}{t}$. Suppose that
      $\dfdx{x}{t}=1\sfrac{\text{ foot}}{\text{second}}$. Is
      $\dfdx{y}{t}$ constant? How does this compare with your guess
      earlier?
    \item  Translate your solution into Newton's dot notation.
    \end{enumerate}
  \end{embeddedproblem-enumerate}
\end{myexample}

\pagebreak[4]
\begin{embeddedproblem-enumerate}{}
  \label{problem:snell1}
    \begin{enumerate}[label={\bf{}(\alph*)}]
    \item   In the following diagram suppose the point $P$ is moving on the
  $x$-axis from left to right with horizontal velocity
  $\dfdx{x}{t} = 1 \sfrac{\text{ unit}}{\text{second}}$.\\
  \centerline{\includegraphics*[height=1.3in,width=2in]{../Figures/PulleyProblem2}}

  Find the distance, $D$ from  $A$ to $P$ to $B$ as a
    single function of $x$ and use this to compute
    $$\eval{\dfdx{D}{t}}{\dfdx{x}{t}}{1}.$$
  \item 
  \begin{enumerate}[label={\bf{}(\roman*)}]
   As an alternative approach, we re-label the diagram as follows:\\
    \centerline{\includegraphics*[height=1.5in,width=2in]{../Figures/PulleyProblem3}}
   \item Define $ D=z+w.$ From  this diagram we have
    \begin{align*}
      w^2\amp{}=9+y^2,\\
      5\amp{}=x+y, \text{ and}\\
      z^2\amp{}= 1+x^2.
    \end{align*}
  \item 
    Differentiate each of these equations and use the results to
    show  that
    $$
    \dx{D}=\left(\frac{x}{z}-\frac{y}{w}\right)\dx{x},
    $$
    then use this to compute $\eval{\dfdx{D}{t}}{x}{1}.$
  \end{enumerate}
\end{enumerate}
\end{embeddedproblem-enumerate}
The purpose of Problems\#\ref{problem:TwoCarts} and~\#\ref{problem:snell1} is to demonstrate that
you have alternatives when solve any problem. You should consciously
try to avoid getting locked in to a single solution scheme.  In this
case one alternative was to try to determine a single function of a
single variable before differentiating.  At the other extreme we can
just assign variable names to all of the important quantities in our
problem and differentiate the equations relating these variables.

Neither approach is necessarily better; that depends upon the
problem. In fact most of the time you'll find yourself using a
mixture of the two. It is always good to have alternatives. Keep this
in mind.


  \begin{embeddedproblem}{}
    Suppose we have a rectangular box whose length is increasing by
    $3$ cm/sec, width is increasing by $2$ cm/sec and height is
    decreasing by $3$ cm/sec.  How fast are the volume and surface
    area of the box changing when the length is $25$ cm, width is $20$
    cm, and height is $10$ cm?  Are they increasing or decreasing?
  \end{embeddedproblem}

%  \pagebreak[4]
\vskip-4mm{}

\begin{embeddedproblem}{}
  Suppose a point $P$ is moving along the curve $y=x^2$ so that its
  horizontal velocity is
  $\dfdx{x}{t}=1\sfrac{\text{ unit}}{\text{second}}$.  For which
  values of $x$ is the distance from $P$ to the point $(0,1)$
  increasing and for what values is it decreasing?  What can you say
  about the points where the distance transitions from increasing to
  decreasing or decreasing to increasing?
\end{embeddedproblem}

% \vskip-2mm{}
% \begin{embeddedproblem}{}
%   The illumination of an object by a light source is inversely
%   proportion to the square of the distance between the source and the
%   object.  Consider the following diagram of two street lamps
%   of equal strength and an object at point P between them. \\
%   \centerline{\includegraphics*[height=1.4in,width=5.5in]{../Figures/StreetlampProblem}}
%   If $P$ is moving away from the left street lamp at a rate of $5
%   \sfrac{\text{ ft}}{\text{second}}$, then how fast is the illumination decreasing when $P$ is
%   $20$ ft from the left street lamp?  When it is $50$ ft from the left
%   street lamp?  Where is the illumination decreasing faster?  Does
%   this \emph{feel} right\aside{That is, does your answer match your
%     intuition?} to you?  Why?  \comment{Use $k$ for the constant of
%     proportionality.}
% \end{embeddedproblem}


\begin{embeddedproblem}{}
    \label{problem:IdealGas}
    The ideal gas law (\href{https://mathshistory.st-andrews.ac.uk/Biographies/Boyle/}{Boyle}'s Law) states that the pressure $P$ (in pascals),
    volume $V$ (in cubic meters), and temperature $T$ (in degrees
    \href{https://mathshistory.st-andrews.ac.uk/Biographies/Thomson/}{Kelvin}) of an ideal gas are related by the formula $PV=kNT$ where
    $N$ is the number of gas molecules and $k$ is the \href{https://mathshistory.st-andrews.ac.uk/Biographies/Boltzmann/}{Boltzmann}
    constant. What this says is that for a fixed amount of gas, if the
    volume is held fixed, then the pressure is proportional to the
    temperature and if the temperature is held fixed, then the
    pressure is inversely proportional to the volume.
    \begin{enumerate}[label={\bf{}(\alph*)}]
      % \begin{multicols}{2}
    \item Assuming that we have an enclosed gas, and that $N$ is held
      constant. Find an equation that relates $\dx{P}$, $\dx{V}$, and
      $\dx{T}$, and use this to find $\dfdx{T}{t}$ in terms of
      $\dfdx{P}{t}$ and $\dfdx{V}{t}$.
    \item
      Suppose we have a piston compressing the gas in a
      cylinder as in the following sketch.\\
      \centerline{\includegraphics*[height=1in,width=1.5in]{../Figures/PistonProblem2}}
      Show that for the pressure to be held fixed $\dfdx{T}{t}$ must
      satisfy
      $$
      \dfdx{T}{t} = \frac{T}{h}\dfdx{h}{t}.
      $$
      \comment{Notice that this is independent of the radius of the
        cylinder.}
    \item Suppose that when $h=.2$ meters the temperature is
      $300^\circ K$.  If the piston is moving to the right at a rate
      of $0.001 \sfrac{\text{meters}}{\text{second}}$, how fast should
      the temperature be decreasing at that instant to maintain
      aconstant pressure? 
      % \end{multicols}
    \end{enumerate}
  \end{embeddedproblem}

\vskip-4mm{}
    \begin{embeddedproblem}{}
      A streetlamp is mounted at the top of a $20$ foot pole.  A $6$
      foot tall person is walking away from the base of the pole at a
      constant speed.
      \begin{enumerate}[label={\bf{}(\alph*)}]
      \item Make a guess: Is the length of this person's shadow
        increasing at a constant rate?
      \item Now solve the problem in part (a).
      \end{enumerate}
    \end{embeddedproblem}

\vskip-4mm{}



%\pagebreak[4]
\begin{embeddedproblem}[Find the Patterns]
%   \begin{wrapfigure}[5]{r}{2in}
% \captionsetup{labelformat=empty}
%     \centerline{\includegraphics*[height=2in,width=2in]{../Figures/RelatedRates-RightTri}}
% \label{fig:RelatedRates-RightTri}
% \end{wrapfigure}
    Suppose the position of a point, $P_x$, on the $x$-axis is given
    by $(x(t),0)$ and the position of another point, $P_y$, on the
    $y$-axis is given by $(0,y(t))$. Let $S(t)$ represent the distance
    between $P_x$ and $P_y$.
    \begin{enumerate}[label={\bf{}(\alph*)}]
    \item
      \begin{enumerate}[label={\bf{}\roman*.}]
        % \begin{multicols}{2}
      \item If $x(t)=t$, and $y(t)=t$ show that
        $\dfdx{S}{t}=\sqrt{2}$.
      \item If $x(t)=2t$, and $y(t)=2t$ show that
        $\dfdx{S}{t}=\sqrt{8}$.
      \item If $x(t)=t$, and $y(t)=3t$ show that
        $\dfdx{S}{t}=\sqrt{10}$.
      \item If $x(t)=5t$, and $y(t)=3t$ show that
        $\dfdx{S}{t}=\sqrt{34}$.
      \item Do you see a pattern yet? If you do write it down and try
        to show that it always holds. If you don't make up some more
        simple examples, solve them to gather more evidence, and
        conjecture a pattern when you believe you have see it. Then show
        that the pattern always holds.
        % \end{multicols}
      \end{enumerate}
    \item
      \begin{enumerate}[label={\bf{}\roman*.}]
      \item Show that if $x(t)=t^{\frac12}$, and $y(t)=t^{\frac12}$
        then $\displaystyle\dfdx{S}{t}=\frac{1}{\sqrt{2t}}$. Does this
        make intuitive sense to you?
      \item Show that if $x(t)=t^{\frac13}$, and $y(t)=t^{\frac13}$,
        then
        $\displaystyle\dfdx{S}{t}=\frac{\sqrt{2}}{3\sqrt[3]{t^2}}$.
      \item Show that if $x(t)=t^{\frac14}$, and $y(t)=t^{\frac14}$,
        then
        $\displaystyle\dfdx{S}{t}=\frac{\sqrt{2}}{4\sqrt[4]{t^3}}$.
      \item Do you see a pattern yet? If you do write it down and try
        to show that it always holds. If you don't make up some more
        simple examples, solve them to gather more evidence, and
        conjecture a pattern when you believe you have one. Then show
        that the pattern always holds.
      \end{enumerate}
    \item Finally, given \emph{any} $x(t)$ and \emph{any} $y(t)$, find
      a completely general formula for $\dfdx{S}{t}$,
    \end{enumerate}
    \vskip-6mm{}
  \end{embeddedproblem}




\section{Self-intersecting Curves and Parametric Equations}
\label{sec:self-inters-curv}
In Problem~\#\ref{problem:ClassicalCurves1} we had you find the line
tangent to several interesting curves.  We were careful however to ask
only about relatively simple curves. In particular none of the curves
in Problem~\#\ref{problem:ClassicalCurves1} intersected themselves. We
avoided self-intersecting curves earlier because we didn't have the
techniques necessary to address the problem of tangent lines at points
of intersection. We do now.

It turns out that we already have all of the computational tools we
need. But we need to broaden the way we think about them.

%\pagebreak[4]
\begin{ProficiencyDrill}
  \label{problem:ClassicalCurves3}
  For the following curves, compute $\dfdx{y}{x}.$
  \begin{enumerate}[label={\bf{}[\Roman*]}]
    \begin{multicols}{2}
    \item Folium of Descartes:\\
      $x^3+y^3=3xy$\\
      \leftline{\includegraphics*[height=1.5in,width=1.5in]{../Figures/FoliumOfDescartes}}
    \item The Tschirnhausen cubic:\\
      $27{y}^{2}=\left(1-x\right){\left(8+x\right)}^{2}$\\
      \leftline{\includegraphics*[height=1.5in,width=3in]{../Figures/TschirnhausenCubic3}}
    \end{multicols}
  \end{enumerate}
  \begin{enumerate}[label={\bf{}(\alph*)}]
    % \begin{multicols}{2}
  \item Show that for the Folium of Descartes
    $\dfdx{y}{x}=\frac{y-x^2}{y^2-x}$. 
  \item    Locate all of the points $(x,y)\neq (0,0)$ where the line tangent to the
    Folium of Descartes is horizontal, and where it is vertical.
  \item Show that for the Tschirnhausen cubic $\dfdx{y}{x} =
    \frac{-(8+x)(3x+6)}{54y}.$
  \item Locate all of the points $(x,y)\neq (-8,0)$ where the line
    tangent to the Tschirnhausen cubic is horizontal, and where it
    is vertical.
    % \end{multicols}
  \end{enumerate}
\end{ProficiencyDrill}



If we try to compute $\eval{\dfdx{y}{x}}{(x,y)}{(0,0)}$ for the Folium
of Descartes we get $$\eval{\dfdx{y}{x}}{(x,y)}{(0,0)}=\frac{0}{0}$$
which is very strange, but seems to make \emph{some} sort of sense
because we can see in the graph that at $(0,0)$ there is both a
horizontal AND a vertical tangent line.

You might be thinking that we have this situation solved; that if we
ever  obtain $\dfdx{y}{x}=\frac{0}{0}$, at some point then the
curve has both a vertical and horizontal tangent at that point.
It would be nice if things were that simple. But they are not.

Next consider at the self-intersection point $(-8,0)$ for the
Tschirnhausen Cubic. Again, as long as we stay away from the
self-intersection point there are no difficulties. You can see from
the graph that any sort of tangent line at the point of
self-intersection would be neither vertical nor horizontal but once
again we have $\eval{\dfdx{y}{x}}{(x,y)}{(-8,0)}=\frac{0}{0}$.

If nothing else, these examples illustrate that using Calculus is much
more subtle than simply computing a derivative. A cavalier attitude
can lead to some very strange anomalies.  As is true of any powerful
tool, to avoid disaster we must be careful.

The safest approach is to avoid points of self-intersection and
fractions such as $\frac00$ but we won't be able to avoid them forever
so we might as well address the issue now.  Adopting Newton's dynamic
approach will give us the crucial insight. We'll need to
think of the horizontal and vertical coordinates as Newton's
fluents; things changing in time.


%\pagebreak[4]
  \begin{embeddedproblem-enumerate}
    \label{problem:TschirnhausenCubic}
    \begin{enumerate}[label={\bf{}(\alph*)}]
      % \begin{multicols}{2}
    \item Show that if $x(t)$ and $y(t)$ satisfy
      \begin{align*}
        x(t)\amp=1-3t^2\\
        y(t)\amp=t(3-t^2)
      \end{align*}
      then $x$ and $y$ satisfy the equation of the Tschirnhausen
      Cubic: $27{y}^{2}=\left(1-x\right){\left(8+x\right)}^{2}$.
    \item Compute the fluxions of $x$ and $y$.
      \item If we think of $t$ as time, with $t\lt 0$, representing
        time in the past, is the point traversing the curve clockwise
        or counterclockwise?
        \\
        \hint{ Consider the values of $t$ for which $x$ and
          $y$ are moving in the positive and negative directions.}
      \item Find the values of $t$ when $(x(t), y(t)) = (-8,0)$ and
        use the result of part (b) to compute
        $\dfdx{y}{x} =\frac{\dot{y}}{\dot{x}}$ at these times.
        Are
        your answers consistent with what you obtained in part (c)?
        Explain.
      \item Show that the equations of the two lines which are tangent
        to the Tschirnhausen cubic at $(-8, 0)$ are
        $
        y=\pm\frac{1}{\sqrt{3}}(x+8).
        $
      % \end{multicols}
    \end{enumerate}
  \end{embeddedproblem-enumerate}


%\pagebreak[4]

\begin{embeddedproblem}
  \label{problem:FoliumOfDescartes1}
  Recall that the Folium of Descartes, $x^3+y^3=3xy$,  crosses itself at the origin.
  % \centerline{\includegraphics*[height=2in,width=2in]{../Figures/FoliumOfDescartes}}
  \begin{enumerate}[label={\bf{}(\alph*)}]
  \item Show that if $ x(t)=\frac{3t}{1+t^3}$ and
    $ y(t)=\frac{3t^2}{1+t^3}$ then $x$ and $y$ satisfy the equation
    of the Folium.
  \item Compute the fluxions of $y$ and $x$.
  \item For what values of $t$ is $y\ge0$?
  \item For what value of $t$ is $(x(t), y(t)) = (0,0)$? Compute
    $\dfdx{y}{x}=\frac{\dot{y}}{\dot{x}}$ at this value of $t$. Is
    this consistent with the graph?
  \item Can you explain why we only found one tangent line when it
    is quite clear from the graph that there must be two? 
  \end{enumerate}
\end{embeddedproblem}

Using the static approach to find a tangent line at a
self-intersection point is very much like standing in the center of
the  intersection of two roads and trying to decide if the road goes
north-south or east-west. Obviously the road doesn't \emph{go}
anywhere. It just sits there. It is the travelers on the road who are
moving.

If instead we ask, ``Which way are we going as we
  \emph{pass through the intersection}?'' there is only one answer. We're
going in the direction we were traveling when we entered the
intersection.  If we think about the situation dynamically rather than
statically then we always have both a position on the curve and a
direction we are traveling.
% Do you see how changing our interpretation of $x$ and $y$ allowed us
% to work with self-intersecting curves?  If we think of $x$ and $y$ as
% coordinates of an unmoving road then it is impossible to interpret
% tangency in any meaningful way. It is like standing in the middle of
% the intersection of two roads and asking, ``Which way does the road
% go?'' Obviously it can be ``going'' (if that is the right word) in any
% of four directions.

A curve is not dynamic. It does not move, it just sits there. Like a
road. But if we think of the curve as being generated by the motion of
a point (Newton's viewpoint) things change.

If we are traveling along a curve it is clear that our direction of
travel is always tangent to the curve. By changing our question from
``What is the tangent line at this \emph{point}?'' to ``What is the
tangent line at this \emph{time}?'' the concept of two (or more)
tangent lines at a \emph{single point in space} becomes
meaningful. Each tangent is obtained by passing through the point of
tangency at \emph{distinct moments in time.}

What comes along with this change in interpretation is this: We no
longer want to think of a differential ratio \emph{only} as a
slope. The ratio $\dfdx{y}{x}$ is simply the rate of change of $y$
with respect to $x$. If $y$ and $x$ happen to represent the vertical
and horizontal coordinates of a graph, then the change of $y$ with
respect to $x$ will be the slope of the curve, $y=y(x)$. If $y$
represents the vertical position of an object and $t$ represents the
time when the object is at position $y(t)$, then the change of $y$ with
respect to $t$, $\dfdx{y}{t}=\dot{y}$, is the vertical velocity of the object
as it passes through $y$.

In general if $\alpha$ and $\beta$ are two related quantities  then
$\dfdx{\alpha}{\beta}$ is the \term{rate of change of $\alpha$ with respect to
  $\beta$.} The physical (or geometric) interpretation of
$\dfdx{\alpha}{\beta}$ will necessarily depend on what $\alpha$ and
$\beta$ represent physically (or geometrically).

Notice that when we switched to the dynamical viewpoint we changed the
nature of our representations of the curves. For example the Folium of
Descartes can be represented by the formula as
$$x^3+y^3=3xy.$$
This formula is complex and can be hard to work with, mainly because
the relationship between $x$ and $y$ is difficult to see and understand.

In the representation
$$
  x(t)=\frac{3t}{1+t^3}, \ \  \  
  y(t)=\frac{3t^2}{1+t^3}
  $$
  the relationship between $x$ and $y$ is still difficult but the
relationship between $x$ and $t$ and between $y$ and $t$ is a bit
simpler because $x$ and $y$ are now both functions of $t$, and we know
how to work with functions.
    
This second representation is usually called the \term{parametric
  functions} or \term{parametric equation} representation, because the
$x$ and $y$ coordinates are given in terms of a parameter. In this
case the parameter $t$ represents time. This is common but by no means
required. The parameter might represent anything, just as $x$ and $y$
might.

In fact, you might be wondering where we got the parameterization for
the Folium.  It may not seem obvious, but as long as $(x,y)\neq(0,0)$, the
parameter $t$ actually represents the slope of the line joining the
origin to $(x,y)$.

\begin{embeddedproblem}
  Let $(x,y)$ represent any point on the Folium $x^3+y^3=3xy$ which is
  not the origin.  Let $t$ represent the slope of the line joining the
  origin to $(x,y)$; that is, $\frac{y}{x}=t$ or $y=tx$.  Use this to
  show that
  $$
  x=\frac{3t}{1+t^3},\text{ and } y=\frac{3t^2}{1+t^3}.
  $$
  \end{embeddedproblem}
  Of course, it is not wrong to also think of this parameter as time.
  We are just stipulating that the point is moving so that its
  position at time $t$ is such that the slope of the line joining the
  origin to the point matches $t$.  Mathematically, what the parameter
  represents is usually not at issue. It is just a parameter.


  \begin{mynotation}{Parametric Functions}
    Recall that when we looked a Roberval's treatment of the conic
    sections in Section~\ref{sec:roberv-conic-sect} we found it handy
    to think of our curves as being traced out by the motion of a
    point and we created the notation
    $
    \begin{ParamEq}
      {x(t)}{y(t)}{}
    \end{ParamEq}
    $ to reflect that point of view. This notation suits our current
    needs almost perfectly, but  we will need to  modify it slightly.

    First, since the coordinates of our point are, individually,
    functions of $t$ it follows that the position of the point $P$
    itself depends on (is a function of) $t$ as well: $P(t)= 
    \begin{ParamEq}
      {x(t)}{y(t)}{}.
    \end{ParamEq}
    $ 

    Second, it is useful for us to have a convenient way to specify
    the domain of this function so we will add a third component to do
    that. For example if the domain of our function is all
    values of $t$ strictly between zero and one we would write.
    $$
    P(t)=
    \begin{ParamEq}
      {x(t)}{y(t)}{0\lt t\lt1}.
    \end{ParamEq}
    $$
  \end{mynotation}

% \pagebreak[4]
% \begin{myexample}\vskip-7mm{}
%   \label{example:CircleParam1}
%   \begin{wrapfigure}[5]{r}{1.5in}
%     \vskip-12mm{}
% \captionsetup{labelformat=empty}
%   \centerline{\includegraphics*[height=1.5in,width=1.5in]{../Figures/CircleParamr}}
% \label{fig:CircleParamr}
% \end{wrapfigure}
%   Suppose $r$ is a constant. Then $
%   P(t)=\begin{ParamEq}
%     {x=r\cos(t)}
%     {y=r\sin(t)}
%     {0\le t \lt 2\pi}
%   \end{ParamEq}
%   $ parameterizes the circle $x^2+y^2=r^2$.  In this
%   example it can be helpful to think of the parameter $t$ as the
%   central angle of the circle as seen below.
%   Since it is customary to represent angles with the Greek letter
%   $\theta$  this parameterization is often rendered as
%   $
%   P(t)=\begin{ParamEq}
%     {x=r\cos(\theta)}
%     {y=r\sin(\theta)}
%     {0\le t \lt 2\pi}
%   \end{ParamEq}
%   $ 
%   but it should be clear the what we call the parameter
%   is unimportant. It is what the parameter represents that matters
%   since this is what guides how we think about it.
% \end{myexample}

``Dynamic'' and
``static'' are only words we use to describe the way we're thinking
about a problem. There is nothing inherently dynamic or static in either representation
so there is no \foreign{a priori} reason to prefer one over the
other. They describe the same set of points in the plane. 

For example if $y=x^2$ then $\dx{y}=2x\dx{x}$. Time ($t$) does not
appear in these formulas so we tend to think of them
statically. However if we want to think of them dynamically we divide
by $\dx{t}$ to get $\dfdx{y}{t}=2x\dfdx{x}{t}$ and interpret this to
say that at any given time $t$, the position of a point is
$
\begin{ParamEq}
  {x(t)}
{y(t)}
{}
\end{ParamEq}
$ and the vertical velocity, $\dfdx{y}{t}$ is twice the
value of the $x$ coordinate times the horizontal velocity,
$\dfdx{x}{t}$.
% It comes down to changing the question from ``What is the tangent line 
% the same way question  then there is only curv in the previous problem, it is hard
% to account for the fact that the Folium crosses itself at the origin.
% We will come back to that later\notetoself{I’m thinking to revisit this in the
% section of limits as  approaches $\pm\infty.$}, but before we go too far afield,
% let’s come back to the problem that prompted us to take Newton’s
% dynamic approach to begin with: the flight path of the Vomit Comet.
% We’ll start by looking at motion in a single direction and then apply
% what we learn to a two dimensional flight path.  To help with the
% jargon, we will introduce some terminology. 


\begin{myexample}
  Moving between the equation and parametric forms can be very hard to
  do depending on the complexity of the equation.  The simplest
  situation is when you have $y$ as a function of $x$; for example
  $y(x)=x^2$. To find a parametric representation we observe that we
  need to specify both $x$ and $y$ as functions of a third parameter,
  $t$. This can be puzzling until we realize that $x$ is completely
  free. All we need to do is ensure that $y=x^2$.  So if we take
  $x(t)=t$ and $y=x^2=t^2$ we \emph{almost} have our
  parameterization.
  
  When faced with a formula like $y=x^2$ you have learned to always
  assume that $x$ could be any real number\aside{At least any real
    number that makes sense.}, but with parametric equations this
  assumption can lead to problems. We'll need to specify the range of
  the parameter $t$ explicitly. This is why we said we \emph{almost}
  have our parameterization. A complete parameterization needs to
  specify the values of $t$ that are available\aside{Otherwise known
    as the domain.} to us, so
  $
P(t)=  \begin{ParamEq}
{    x(t)=  t}
{    y(t)=  t^2}
{    -\infty\lt{}t\lt{}\infty}
  \end{ParamEq}
  $.

  We can \emph{always} parameterize the graph of a function the same
  way we just parameterized $y=x^2$. A parameterization of $y=y(x)$,
  with domain $a\le x\le b$ is
$
P(t)=
\begin{ParamEq}
{  x(t)=t}
{  y(t)=y(t)}
{  t\in [a,b]}
\end{ParamEq}.
$

This parameterization may or may not be particularly useful, but
it can \emph{always} be done. 
\end{myexample}



\begin{embeddedproblem-enumerate}
    \begin{enumerate}[label={\bf{}(\alph*)}]
    \item Show that the parameterization
      $P(t)=
      \begin{ParamEq}
        {t^2+1}
        {t^4+2t^2+1}
        {-\infty\lt{}t\lt{}\infty}
      \end{ParamEq}
      $
      traces out part of the parabola $y=x^2$. Which part?
    \item The parameterization
      $P(t)=
      \begin{ParamEq}
        {t^2+1}
        {t^4+2t^2+1}
        {0\lt{}t\lt{}\infty}
      \end{ParamEq}
      $
      traces out  different part of the same parabola. How else is this
      parameterization different from the one in part (a)?
    \end{enumerate}
  \end{embeddedproblem-enumerate}

  \vskip-3mm{}
\begin{embeddedproblem-enumerate}
  \begin{enumerate}[label={\bf{}(\alph*)}]
  \item   Show that each of the following is a parameterization of part of the graph
    of $y=x^2$.
    \begin{enumerate}[label={\bf{}(\roman*)}]
      \begin{multicols}{3}
      \item $P(t)=
        \begin{ParamEq}
          {3t}
          {9t^2}
          {t\in\RR}
        \end{ParamEq}
        $
      \item $P(t)=
\begin{ParamEq}
  {-t}
  {t^2}
  {t\in\RR}
\end{ParamEq}
$
      \item $P(t)=
\begin{ParamEq}
  {\sqrt{t}}
  {t}
  {t\ge0}
\end{ParamEq}
$
      \item $P(t)=
\begin{ParamEq}
  {-\sqrt{t}}
  {t}
  {t\ge{}0}
\end{ParamEq}
        $
      \item $P(t)=
\begin{ParamEq}
{1/t}
{1/t^2}
{\abs{t}\gt{}0}
\end{ParamEq}
        $
      \item $P(t)=
\begin{ParamEq}
{t^{1/3}}
{t^{2/3}}
{t\gt 0}
\end{ParamEq}
        $
      \item $P(t)=
\begin{ParamEq}
{t^3-1}
{t^6-2t^3+1}
{-\infty\lt{}t\lt{}\infty}
\end{ParamEq}
        $
      \item $P(t)=
\begin{ParamEq}
  {t+1}
  {t^2+2t+1}
  {-20\lt{}t\lt{}20}
\end{ParamEq}
        $
        % \item $
        %   \begin{Bmatrix}
        %     x(t)=\sin(t)\\
        %     y(t)=\sin^2(t)\\
        %     -\infty\lt{}t\lt{}\infty
        %   \end{Bmatrix} $
        % \item $
        %   \begin{Bmatrix}
        %     x(t)=\cos(t)\\
        %     y(t)=1-\sin^2(t)\\
        %     -\infty\lt{}t\lt{}\infty
        %   \end{Bmatrix} $
      \end{multicols}
    \end{enumerate}
  \item Sketch \emph{only} the part of the curve included in each
    parameterization in part (a).  Be sure to indicate the direction
    of travel in each case, assuming $t$ is increasing.
  \item Compute $\frac{\dot{y}}{\dot{x}}$ and show that this yields
    $\dfdx{y}{x}=2x$ for each of the parameterizations in part (a).
  \end{enumerate}
\end{embeddedproblem-enumerate}

  \vskip-3mm{}
\begin{embeddedproblem-enumerate}
  \label{problem:UnitCircleParameterization1}
  \begin{enumerate}[label={\bf{}(\alph*)}]
  \item Show that each of the following is a parameterization of
    part of the unit circle.
    \begin{enumerate}[label={\bf{}(\roman*)}]
      \begin{multicols}{3}
      \item $P(t)=
        \begin{ParamEq}
{          t}
{          \sqrt{1-t^2}}
{          -1\le t\lt{} 1}
        \end{ParamEq}
        $
      \item $P(t)=
        \begin{ParamEq}
{          t}
{          -\sqrt{1-t^2}}
{          0\le t\lt{} 1}
        \end{ParamEq}
        $
      \item $P(t)=
        \begin{ParamEq}
{          \sqrt{1-t^2}}
{          t}
{          -1\le t\lt{} 1}
        \end{ParamEq}
        $
      \item $P(t)=
        \begin{ParamEq}
{          \sqrt{1-t^2}}
{          t}
{          0\le t\lt{} 1}
        \end{ParamEq}
        $
      \item $P(t)=
        \begin{ParamEq}
{          t^2}
{          \sqrt{1-t^4}}
{          -1\le t\lt{} 1}
        \end{ParamEq}
        $
      \item $P(t)=
        \begin{ParamEq}
{          t-2}
{          \sqrt{4t-(3+t^2)}}
{          1\le t\lt{} 2}
        \end{ParamEq}
        $
      \end{multicols}
    \end{enumerate}
  \item Sketch \emph{only} the part of the curve included in each
    parameterization in part (a).  Be sure to indicate the direction
    of travel in each case.
  \item Compute $\frac{\dot{y}}{\dot{x}}$ and show that this yields
    $\dfdx{y}{x}=-\frac{x}{y}$ for each of the
    parameterizations in part (a).
  \end{enumerate}
\end{embeddedproblem-enumerate}

As we've seen we can always parameterize the graph of a function, but
the reverse is not true.  A parameterized curve will not always be the
graph of some function.  For example none of the curves in
problem~\#\ref{problem:ClassicalCurves1} is the graph of a function,
but all can be parameterized.

Since it can't always be accomplished there is no general strategy for
expressing a parameterized curve as the graph of a function. One way
to do it\aside{Not necessarily the best.  In general there is no best
  way.} is to find $\dfdx{y}{x}$ and ``undifferentiate''  as in the
following example.

%\pagebreak[4]
\begin{embeddedproblem}
  \label{problem:ParametricIVP}
  \vskip-4mm{}
  \begin{wrapfigure}[]{r}{2in}
  \vskip-1mm{}
    \captionsetup{labelformat=empty}
    \centerline{\includegraphics*[height=2in,width=2in]{../Figures/Parametric4}}
    \label{fig:Parametric4}
  \end{wrapfigure}
  The red curve  in the sketch at the right 
  is  parameterized by
  $$P(t)=
  \begin{Bmatrix}
    \frac{1}{\sqrt{t^2-1}}\\[6pt]
    \frac{t^2-2\sqrt{t^2-1}}{\sqrt{t^2-1}}\\[6pt]
    1\lt{} t\lt{}\infty
  \end{Bmatrix}
  .
  $$
  We wish to find
  $y(x)$.
  Differentiating $x$ and $y$ gives
  $$\dfdx{x}{t}= \frac{-t}{(t^2-1)^{\frac32}}$$
  and
$$
    \dfdx{y}{t}= \frac{t^3-2t}{(t^2-1)^{\frac32}}
$$
so that
$
\dfdx{y}{x}
% \frac{\dfdx{y}{t}}{\dfdx{x}{t}}\frac{\frac{t^3-2t}{(t^2-1)^{\frac32}}}{\frac{t}{(t^2-1)^{\frac32}}}=
% -t^2+2
= 1-(t^2-1)$.
  % and since $x(t)=\frac{1}{\sqrt{t^2-1}}$ we have
  % $\dfdx{y}{x} =1-\frac{1}{x^2},$ right?
  \begin{enumerate}[label={\bf{}(\alph*)}]
%    \begin{multicols}{2}
    \item    Fill in the details of the computation of $\dfdx{x}{t}$,
      $\dfdx{y}{t}$,
    and show that  $\dfdx{y}{x} = 1-\frac{1}{x^2}$.
  \item Show that  $\dfdx{y}{x} = 1-\frac{1}{x^2}$ is the derivative
    of $y=x+\frac{1}{x}$.
  \item Sketch the graph of $y= x+\frac1x$ to see that this is \alert{not} the
    same as the parameterized curve above.  Can you tell what went
    wrong?
    \item Show that the point $\left(1,0\right)$ is on our
      parameterized  curve, and use this to find
      the correct function, $y(x)$.
  % \end{multicols}
  \end{enumerate}
\end{embeddedproblem}  

% \begin{embeddedproblem-enumerate}
%     \begin{enumerate}[label={\bf{}(\alph*)}]
%     \end{enumerate}
%   \end{embeddedproblem-enumerate}

\pagebreak[4]

\begin{embeddedproblem}
  For each of the given parameterizations find $\dfdx{y}{x}$ two
  different ways:
  \begin{enumerate}[label={\bf{}(\roman*)}]{}
  \item First by computing $\dx{x}$ and $\dx{y}$ and finding
    their ratio. 
  \item and second by solving for $t$ in terms of $x$ (or $y$),
    substituting this into $y$ (or $x$) and computing $\dfdx{y}{x}$
    directly.
  \end{enumerate}

  Then use this to determine the corresponding equation in $x$ and
  $y$.
  \begin{enumerate}[label={\bf{}(\alph*)}]
    \begin{multicols}{3}
    \item $P(t)=
      \begin{ParamEq}
        {t-1}
        {t+1}
        {t\in\RR}
      \end{ParamEq}
      $
    \item $P(t)=
      \begin{ParamEq}
{        3t-2}
{        -5t+7}
{        t\in\RR}
      \end{ParamEq}
      $
    \item $ P(t)=
      \begin{ParamEq}
{        1/\sqrt{t}}
{        \sqrt{t}}
{        0\lt{}t}
      \end{ParamEq}
      $
    \end{multicols}
  \end{enumerate}
\vskip-10mm{}
\end{embeddedproblem}

% \begin{embeddedproblem}
%   Suppose that $
%   \begin{Bmatrix}
%     x(t)=z(t)\\
%     y(t)=w(t)\\
%     a\le t\lt{} b
%   \end{Bmatrix} $ is a parameterization of some curve. Show that
%   $ \dfdx{y}{x}=\frac{\dfdx{w}{t}}{\dfdx{z}{t}}.  $
% \end{embeddedproblem}



% \begin{embeddedproblem}{}
%   \label{problem:cycloid1}
%   Suppose we trace the path of a point $P$ on a unit circle as the circle
%   rolls along a straight line.  Such a curve is called a \href{https://mathshistory.st-andrews.ac.uk/Curves/Cycloid/}{\term{cycloid}}.\\
%   \centerline{\includegraphics*[height=2in,width=5in]{../Figures/Cycloid1}}
%   For reasons we will examine later this curve has fascinated
%   mathematicians for a long time. For now we just want to find a
%   parameterization for it.
%   Use the above diagram to show that the coordinates of $P$
%   are given by $
%   \begin{Bmatrix}
%     x=\theta-\sin(\theta) \\
%     y=1-\cos(\theta)\\
%     -\infty\lt{}\theta\lt{}\infty
%   \end{Bmatrix}.
%   $
%   \\    % \begin{enumerate}[label={\bf{}(\alph*)}]{}
%     % \item Use the above diagram to show that the coordinates of $P$
%     %   are given by
%     %   \begin{align*}
%           %     x\amp{}=a(\theta-\sin(\theta)) \\
%           %     y\amp{}=a(1-\cos(\theta))
%                          %   \end{align*}
%                          % \item Compute $\dx{x}$ and $\dx{y}$ in terms of $\dx{\theta}.$
%                          % \item Find the slope of the tangent line to the cycloid at $P$ in
%                          %   terms of $\theta$ and use this to confirm that the tangent line is
%                          %   horizontal at the apex of the cycloid.
%                          % \item Suppose the ball is rolling at a speed of
%                          %   $2\frac{\text{units}}{\text{second}}.$ Compute the speed,
%                          %   $\dfdx{s}{t}= \frac{\sqrt{(\dx{x})^2+(\dx{y})^2}}{\dx{t}}$ of
%                          %   the point $P.$ What do you notice about the speed of $P$ in
%                          %   relation to the ball's radius? Do you find this surprising? When
%                          %   is $P$ moving the fastest and when is it moving the slowest?
%                          %   Does this make sense physically?
%                          % \end{enumerate}
% \end{embeddedproblem}






% \vskip1in{}
\section{Bridges, Chains, Domes, and Telescopes}
\label{sec:bridg-chains-telesc}
\markright{\sc{}Bridges, Chains, Domes, and Telescopes}

\subsection*{\underline{Bridges}}
In the previous section all of the quantities we were looking at were
time dependent -- they were in motion -- so it was useful to think of
the derivative, $\dfdx{y}{t}$, as a velocity (fluxion). In other
applications, where the quantities involved are not time dependent,
the derivative is still a useful tool but thinking of it as a velocity
may not be as useful.
\vskip1mm{}
\centerline{\includegraphics*[height=2.3in,width=4in]{../Figures/GoldenGatebyModestosUrbanos1.png}}
\vskip1mm{}


For example, consider the shape of the support cable on a stable
suspension bridge as seen in the picture above. A \emph{stable} bridge
is one that does not move in any significant way\aside{The whole point of
  a bridge is that it does \emph{not} move. A moving bridge is called a
  disaster.}.

% \centerline{\includegraphics*[height=2in,width=4in]{../Figures/MackinacBridge1}}
The shape of the suspension cable appears to be parabolic, but how can
we be sure of this?

At any given moment there are a lot of forces acting on a bridge, but
most of them are insignificant most of the time. The mantra of
mathematical modeling is to keep things as simple as possible (at
least in the beginning) by focusing on the dominant parameter
first\aside{Or, equivalently, we set the secondary parameters equal to
  zero.} In this case the dominant force involved is clearly
the weight of the bridge itself.

We assume that the weight density of the deck of the bridge is $W$
newtons/meter\aside{A \emph{newton} is a unit of force obviously named
  after Isaac Newton.}. The cable adds to the weight of the bridge as
well, but since it is relatively small compared to the weight of the
deck we will ignore it (to keep the model simple).  Assuming the deck
is horizontal we can take it to be the $x$ axis Now suppose we graph
of $y=y(x)$ represents the cable between two upright towers located at
$x=-U$ and $x=U.$ Remember, we don't yet know if this
is a parabola.\\\\
\centerline{\includegraphics*[height=1in,width=5in]{../Figures/Bridge1}}
Let's choose an arbitrary point $(x,y)$ on the curve with $0\lt{}x\lt{}U$ and
focus on the portion of the bridge between $(-x,y)$ and $(x,y)$.\\\\
\centerline{\includegraphics*[height=1in,width=5in]{../Figures/Bridge2}}
% \begin{wrapfigure}[]{r}{1in}
%   % \captionsetup{labelformat=empty}
%   \centerline{\includegraphics*[height=2in,width=1in]{../Figures/VerticalCable1}}
%   \label{fig:VerticalCable1}
% \end{wrapfigure}
These are the two points that are holding up the section of the bridge
on the interval from $-x$ to $x$.  The arrows represent the tangential
forces felt on the cable at those two points.  We assume that the
curve is symmetric so we can restrict our attention to the right half
of the bridge.

Let’s draw that half and examine the tangential force at various
points on the bridge and the horizontal and vertical components of
these.\\
\centerline{\includegraphics*[height=1.7in,width=6in]{../Figures/Bridge6}}
Notice that at $x=0$, the tangential force, which we labeled $-T$ is
horizontal. You’ll see why we used negatives momentarily. This
reflects the fact that at that point there is no vertical force as it
would be the weight for the deck from $x=0$ to $x=0$ which is
zero. (If you think about it, this is why it’s the lowest point, there
is no vertical pull.) Its horizontal component, $-H$, represents the
horizontal tension. Basically, this represents the force opposing us
as we pull the cable tighter. The tighter we winch the cable, the
larger $H$ would be. Notice that this horizontal tension is the same at
every single point since tightening the cable produces the same
horizontal tension throughout the cable. What varies is the vertical
component of the force. As we move out further along the cable, there
is more of the deck to support and since we are ignoring the weight of
the cable, this vertical force is just the weight of the deck from
$x=0$
to that point.

This cumulative effect can also be thought of this way. The point on
the cable at $x = 20$ not only must support the weight of the deck
from $x = 10$ to $x = 20$, but must bear the brunt of the load at the
point $x = 10$ as well. This leads to the vertical force at $x = 20$
being $-10W - 10W = -20W$. The point at $x = 30$ must support whatever
weight there is at $x = 20$ along with the extra weight from $x = 20$
to $x = 30$ which totals $-20W - 10W = -30W$. If you were hoisting a
person carrying a bucket of sand, then you would be bearing the weight
of both the person and the sand whereas that person is only bearing
the weight of the sand. This is the same idea.

The conclusion of this analysis is that at any point
$(x, y)$ on the curve, the tangential force pulling against us has a
horizontal component of $-H$ and a vertical component of $-Wx$. To
stabilize the bridge, we need to be pulling with the same force in
exactly the opposite direction. This leads to the following
drawing. (Now you see why we used negatives before.)\\
\centerline{\includegraphics*[height=1.7in,width=5.8in]{../Figures/Bridge7}}

\begin{embeddedproblem-enumerate}
  \begin{enumerate}[label={\bf{}(\alph*)}]
  \item Use the above analysis and schematic drawing of the forces
    involved to show that the curve which represents the cable must
    satisfy the differential equation
    $$\dfdx{y}{x}=\frac{Wx}{H}.$$
    \hint{Notice that $T$ is pulling in the direction tangent to the
      cable.}
  \item Show that the parabola $y=\frac{W}{2H} x^2+b,$ where $b$ is an
    (unknown) constant, satisfies this differential equation
  \item Let $W = 1$, $b = 0$ and plot the parabola for $H = 1, 2,
    3$. Does this agree with the idea that $H$ indicates how tight we
    are winching the cable?
  \end{enumerate}
  \vskip-4mm{}
\end{embeddedproblem-enumerate}

With our simplifying assumptions in place it appears that  the cable on a
suspension bridge does indeed  hang in the shape of a parabola. But the
suspension bridge has to hold up the roadway beneath it. After all,
that is its purpose.

What about a hanging chain that only has to hold  its own weight?

\subsection*{\underline{Chains}}

What shape  do you think a chain takes on if it is pinned at both
ends, and allowed to hang freely in between? Take
  your best guess and don’t worry about guessing correctly. We'll
  solve this problem shortly.


If you guessed that the curve is a parabola, then you are in good
company.  Even Galileo  believed that a hanging chain assumed
a parabolic shape as we see in  the following passage from his
\href{https://www.maa.org/press/periodicals/convergence/mathematical-treasure-galileos-two-new-sciences}{Dialogue Concerning Two New Sciences} ($1638$).

  \begin{minipage}[c]{.85\linewidth}
    \begin{wrapfigure}[]{r}[.5in]{2in}
      \captionsetup{labelformat=empty}
      \vskip-4mm{}
      \centerline{\includegraphics*[height=1.85in,width=2.1in]{../Figures/ChainWithoutBob}}
      \label{fig:ChainWithoutBob}
    \end{wrapfigure}
    \it{} ``Drive two nails into a wall at a convenient height and at
    the same level; make the distance between these nails twice the
    width of the rectangle upon which it is desired to trace the
    semiparabola. Over these two nails hang a light chain of such
    length that the depth of its sag (curve or \foreign{sacca}) is
    equal to the length of the prism. This chain will assume the form
    of a parabola, so that if this form be marked by points on the
    wall we shall have described a complete parabola which can be
    divided into two equal parts by drawing a vertical line through a
    point midway between the two nails .\ .\ .\ Any ordinary mechanic
    will know how to do it.''
  \end{minipage}

\vskip2mm{}
  \begin{embeddedproblem}[The Hanging Chain]
    \vskip-4mm{}
    \begin{wrapfigure}[10]{r}{2.2in}
      \captionsetup{labelformat=empty}
      \centerline{\includegraphics*[height=1.4in,width=2.2in]{../Figures/HangingChainSchematic}}
      \label{fig:HangingChainSchematic}
    \end{wrapfigure}
    This problem is very similar to finding the shape of the
    suspension bridge cable.  The difference here is that since there
    is no deck, the only vertical force will be the weight of the
    chain itself, so this time the weight of the chain can't be
    ignored.

    Let $w$ represent the weight density of the chain in newtons per
    meter and $s$ represent the length of the chain from the lowest
    point to $P$ and show that the curve represented by the chain must
    satisfy the differential equation
    \begin{equation}
      \dfdx{y}{x}=\frac{ws}{H}\label{eq:Catenary1}
    \end{equation}
    \vskip-.2in{}
  \end{embeddedproblem}
  
  OK, but so what? All equation~(\ref{eq:Catenary1}) says is that the
  proportion $\dfdx{y}{x}$ is equal to the proportion $\frac{ws}{H}$.
  
  If we could \emph{solve} this equation for $y(x)$ then the graph of
  $y$ would be the shape of the chain, but sadly we can't do that
  (yet).

  So why did we bother writing down
  equation~(\ref{eq:Catenary1}) at all? Does this tell us anything
  about the shape of the hanging chain?
  
  Yes. Actually it does.

  The fact is that in our attempts to describe the world around us
  using mathematics we can't actually solve \emph{most} of the
  equations we write down. Real world phenomena are just too
  complex. But just writing them down is a step forward.  For example,
  although we can't (yet) solve equation~(\ref{eq:Catenary1}), and
  thereby find the shape of a hanging chain, we can use it to
  eliminate some shapes.  In particular we can now show that Galileo
  was wrong about the shape of a hanging chain.


\begin{embeddedproblem}[Galileo was wrong!]
  \label{PIC:catenary} 
  \vskip-.3in{}
  The difficulty is that we have the wrong variable on the right side
  of equation~\ref{eq:Catenary1}. Since $y$ depends on $x$ we'd like
  to have only the variables $x$ and $y$, appearing in our
  equation. Instead we have $x$, $y$, and $s$. And we really don't
  know anything about $s$.

  However, recall that we do know something about $\dx{s}$.

  \begin{enumerate}[label={\bf{}(\alph*)}]{}
  \item Use a differential triangle to show that the hanging chain
    curve must satisfy the differential equation
    \begin{equation}
    \dfdxn{y}{x}{2}=\frac{w}{H}  \dfdx{s}{x}=\frac{w}{H}
    \sqrt{1+\left(\dfdx{y}{x}\right)^2 }.\label{eq:CatenaryDiffeq}
  \end{equation}
  \hint{Notice that $\dx{s}$ does not appear in this problem, only
    $s$. Find a way to bring $\dx{s}$ into play starting with
    equation~\ref{eq:Catenary1}.}
  \item Show that the general parabola $y=ax^2+bx+c$ does \underline{not} satisfy this
    differential equation (and so Galileo was mistaken!)  .
  \item Show that, in fact no nonzero polynomial
    $y=a_0+a_1x+a_2x^2+\ldots+a_nx^n$ satisfies the equation of the
    catenary. \\
    \hint{Suppose it did satisfy the equation.  If we were
      to square both sides of the equation, what would be the degree of
      each side?}
  \end{enumerate}
  \vskip-.3in{}
\end{embeddedproblem}



Galileo was a brilliant man and his brilliance is not at all
diminished by this error.  As we said in
Chapter~\ref{cha:problem-solving}, \alert{Smart people make mistakes!
  This is how they get to be smart.} You should follow the example of
Galileo (and Newton, and Leibniz) by making and embracing lots of
mistakes. But don't forget to learn from them as well.

Also, in fairness to Galileo, he did not have access to  the Calculus tools that
make this problem tractable. We do.


\subsection*{\underline{Domes}}
The previous problem shows that the shape of a hanging chain is not a
parabola as Galileo believed. But it actually shows much more than
that. It shows that the shape of a hanging chain is not the graph of
\alert{any} polynomial regardless of its degree. (A parabola is the
graph of a polynomial of degree $2$.) This says that the set of
polynomials is too small to model all of the phenomena we observe in
the world. To solve this problem we'll need a larger class of
functions.

The curve which \emph{does} satisfy
equation(~\ref{eq:Catenary1}), and thus is the shape of a hanging
chain is called the catenary. This is
not especially illuminating since
the word ``catenary'' is derived from the Latin word {\it{}catena}
which simply means chain.



\begin{wrapfigure}[14]{r}{2.7in}
%  \vskip-6mm{}
  \captionsetup{labelformat=empty}
  % \href{https://uh.edu/engines/epi1169.htm}
  % {
  % \href{https://uh.edu/engines/epi1169.htm}{\includegraphics*[height=2.1in,width=3.3in]{../Figures/Pantheon}}
  \href{https://uh.edu/engines/epi1169.htm}{\includegraphics*[height=1.7in,width=2.7in]{../Figures/Pantheon}}
  % }
  \caption{\textcolor{black}{{Drawings of the Pantheon and a cross section of the dome with an inverted catenary superimposed on the arch.} }}
  \label{fig:Pantheon}
\end{wrapfigure}
The last line in Galileo's quote (above) is interesting. The structural
properties of the catenary have been known and used since
ancient times

The Pantheon in Rome was completed during the reign of Hadrian and
dedicated around $126$~AD. It is not known who designed it, but there
is evidence\aside{For example, its shape.} which suggests that a
catenary design was employed.

The
architect
\href{https://www.britannica.com/biography/Filippo-Brunelleschi}{Filippo
  Brunelleschi} ($1377-1446$) used a hanging chain to develop the
shape of the dome of the
\href{http://www.museumsinflorence.com/musei/cathedral_of_florence.html}{Cathedral
  of Santa Maria del Fiore}{} in Florence, Italy in the Middle Ages.

A number of modern architects and masons have used the catenary to design
arches and domes (a rotated
arch) as well.


 For example, the mathematician and architect
 \href{https://mathshistory.st-andrews.ac.uk/Biographies/Wren/}{Sir
   Christopher Wren} ($1632-1723$) used the shape of the catenary arch
 in the  design
 and construction of the structural part of the dome on
 \href{https://en.wikipedia.org/wiki/St_Paul%27s_Cathedral}{St Paul's
   Cathedral in London.}

 While designing St. Paul's Wren was advised by his good friend
 \href{https://mathshistory.st-andrews.ac.uk/Biographies/Hooke/}{Robert
   Hooke} ($1635-1703$), who once made the following  observation about the
 catenary\aside{Hooke is also famous for showing that the contracting
   force of a stretched spring is proportional to the amount of
   stretching. This is known as ``Hooke's Law.'' We will examine this
   in more depth in Section~\ref{sec:simple-harm-oscill}. }
 \begin{center}
   \begin{minipage}[]{.7\linewidth}
     ``As hangs a flexible cable so, inverted, stand the touching
     pieces of an arch.''
   \end{minipage}
 \end{center}
  
% \\\\
  % \centerline{
  %   \includegraphics*[height=2.8in,width=5.2in]{../Figures/StPauls}}
  % \textcolor{blue}{\centerline{ {\it{}St. Paul's Cathedral, London,
  %       UK}}}
 
  \begin{wrapfigure}[10]{l}{2.8in}
    \vskip-3mm{}
    \captionsetup{labelformat=empty}
%    \href{https://en.wikipedia.org/wiki/St_Paul%27s_Cathedral}{\includegraphics*[height=1.5in,width=3.5in]{../Figures/StPauls}}
    \href{https://en.wikipedia.org/wiki/St_Paul%27s_Cathedral}{\includegraphics*[height=1.2in,width=2.8in]{../Figures/StPauls}}
    \label{fig:StPauls}
    \textcolor{blue}{\centerline{ {\it{}St. Paul's Cathedral,}}}
          \textcolor{blue}{\centerline{ {\it{}London, UK}}}
 \end{wrapfigure}
 The previous problems actually indicate why this structure is so
 stable.   The idea to use catenaries in building design was most likely
  originally motivated by experimentation and intuition. But after
  Calculus was invented it became possible to explain why they are so
  stable.
 When the forces involved in causing a hanging chain to fall
 into the shape of a catenary are inverted (turned upside down) then
 the weight of the catenary arch is directed tangentially along the
 arch toward the base.  Catenary arches are still in use today in part
 because they have this property.

  % A contemporary of Newton and Leibniz, The mathematician and
  % architect
  % \href{https://mathshistory.st-andrews.ac.uk/Biographies/Wren/}{Sir
  %   Christopher Wren} ($1632-1723$) used this property in the design
  % and construction of
  %  the structural part of the dome on St Paul's
  % Cathedral in London.



% \begin{wrapfigure}[6]{r}{1.4in}
%   \vskip-8mm{}
%   \captionsetup{labelformat=empty}
%   % \href{http://www.museumsinflorence.com/musei/cathedral_of_florence.html}{\includegraphics*[height=1.7in,width=3in]{../Figures/SantaMariaDelFiore}}
%   \href{http://www.museumsinflorence.com/musei/cathedral_of_florence.html}{\includegraphics*[height=.75in,width=1.5in]{../Figures/SantaMariaDelFiore}}
%   \label{fig:SantaMariaDelFiore}
%    \textcolor{blue}{\centerline{ {\it{}Cathedral
%          of}}}
%    \textcolor{blue}{\centerline{ {\it{}Santa Maria del Fiore}}}
% %     \textcolor{blue}{\centerline{\it{}Florence, Italy}}
% \end{wrapfigure}

  \begin{wrapfigure}[5]{r}{1.1in}
    \vskip-8mm{}
    \captionsetup{labelformat=empty}
    \centerline{\href{https://commons.wikimedia.org/wiki/File:Budapest\_Keleti\_teto\_1.jpg}{
        % \includegraphics*[height=1.75in,width=2in]{../Figures/EastRailStation}}}
        \includegraphics*[height=1in,width=1.1in]{../Figures/EastRailStation}}}
    \label{fig:EastRailStation}
    \caption{\centerline{\textcolor{black}{East Rail
          Station}}\\\centerline{\textcolor{blue}{Budapest, Hungary}}}
        \end{wrapfigure}
        Because of its inherent strength and stability the catenary
        arch is still a common choice for arches and domes today. For
        example, the roof of the East Rail Station in Budapest, Hungary
        is clearly also a catenary.


  We've
  seen that a catenary is not a parabola, but we still don't have a
  formula for the catenary curve.

  So far.

We will return to, and answer, this  question in
Problem~\#\ref{problem:HangingChain}.

  \subsection*{\underline{Telescopes}}
  All telescopes follow one of two basic designs, refractive
  and reflective. Both kinds magnify images by focusing light but they
  do it in two entirely different ways. Refractive telescopes use the
  refractive properties of light\aside{Hence the name.} described by
  Snell's Law (see Section~\ref{sec:snells-law-refr}) to focus the
  light rays of an image. Reflective scopes use the reflective
  properties\aside{Hence the name.} of the
  conic sections we studied in Chapter~\ref{chapter:ScienceBeforeCalc}
  to focus the light rays of an image.

    \begin{wrapfigure}[12]{l}{2.5in}
    \vskip-4mm{}
    \captionsetup{labelformat=empty}
    \href{https://www.lbto.org/}{\includegraphics*[height=1.5in,width=2.5in]{../Figures/MountGraham}}
    % \centerline{\it{}The Large Binocular Telescope at Mt. Graham, Arizona}
    \caption{\textcolor{black}{{\centerline{The Large Binocular Telescope}}\\ \centerline{Mt. Graham, Arizona}}}
    \label{fig:MountGraham}
  \end{wrapfigure}
  Galileo was the first person to point his (refractive) telescope at
  the sky and write about what he saw.  A refractive telescope uses
  two lenses to magnify
an image by refracting light toward the
  observer's eyepiece. But refraction separates the  frequencies
  (colors) of light so refractive telescopes tend to cause chromatic
  aberrations in the image\aside{Little rainbows appear at the edges
    of the objects in view.}

  Because of his work in optics, Newton realized that he could use the
  reflective properties of the parabola to achieve greater
  magnification with a smaller physical device while simultaneously
  avoiding the chromatic aberrations.  All modern telescopes are built
  on Newton's original design which uses a reflective, parabolic
  (primary) mirrors to reflect light toward the focus of the
  par\-a\-bo\-la (remember Roberval) where the flat (secondary) mirror
  reflects it to the eyepiece.  Large telescopes tend to be reflective
  telescopes as they really don't need the tube (which actually only
  holds the two lenses in the refractive telescope in place).  As an
  example, consider the Large Binocular Telescope at Mt. Graham,
  Arizona, seen below.

  % \begin{wrapfigure}[12]{r}{4.2in}
  %   \vskip-3mm{}
  %   \captionsetup{labelformat=empty}
  %   \centerline{\includegraphics*[height=1.8in,width=4.2in]{../Figures/SpinCasting7}}
  %   \label{fig:SpinCasting7}
  % \end{wrapfigure}

  \begin{wrapfigure}[11]{r}{3.8in}
    \vskip-7mm{}
    % \captionsetup{labelformat=empty}
    \centerline{\includegraphics*[height=1.8in,width=3.8in]{../Figures/SpinCasting1}}
    \label{fig:SpinCasting1}
  \end{wrapfigure}
  The Richard F. Caris Mirror Laboratory uses a process called spin
  casting to manufacture these large mirrors. Each mirror is $8.4$
  meters in diameter.  The image at the right gives you an idea of the
  size of one of these mirrors.% \\
  % \centerline{\includegraphics*[height=1.8in,width=3.8in]{../Figures/SpinCasting1}}

  How does one make such a large,
  perfectly parabolic mirror?  To spin cast such a large mirror, the
  people at the Caris Laboratory place high quality borosilicate glass
  into a revolving furnace.  The furnace spins as the glass
  liquefies.
  
  Obviously, when the molten glass spins the middle goes down and the
  sides go up.  What is not obvious is that the surface generated must
  be a parabola.  Clearly, the shape of the mirror, and hence its focal
  length, will be affected by the rotational speed of the glass. The
  question is, how fast must we spin the furnace to produce a
  particular focal length?

  To begin to answer this we will need to think about a typical point
  mass at a point on the surface of the molten glass as shown below and
  suppose the glass is spinning at an angular
    \begin{wrapfigure}[11]{l}{3in}
%    \vskip-4mm{}
\captionsetup{labelformat=empty}
  \centerline{\includegraphics*[height=1.7in,width=3in]{../Figures/Spincasting8}}
\label{fig:Spincasting8}
\end{wrapfigure}
velocity of
  $\omega$~$\sfrac{\text{radians}}{\text{second}}$.  Our task is to find
  the function $y=y(x)$ whose graph gives the shape of the surface of
  the liquid glass.
  

The upward pointing arrow at the left depicts the force which keeps
the point mass elevated and spinning in a circle. It will be
perpendicular to the surface of the liquid. (Think about a hose with a
hole in it.  The water sprays out in a stream perpendicular to the
hose.)

\begin{wrapfigure}[]{r}{1in}
\captionsetup{labelformat=empty}
      \centerline{\includegraphics*[height=.9in,width=1in]{../Figures/SpinCasting4}}
\label{fig:SpinCasting4}
\end{wrapfigure}
If we separate that force into its vertical and horizontal components,
  then we get the diagram at the right.  The magnitude of vertical
  component of the force is $mg$, where $g$ is the acceleration due to
  gravity.  This is the force needed to counter the weight of the
  particle.  The only horizontal force is the centripetal force which is
  keeping the particle moving in a circle around the central axis. In
  the next section we will see that the magnitude of the centripetal
  force is $mx\omega^2$, so for now we will assume this is correct and
  proceed.

  \begin{embeddedproblem-enumerate}
    \label{PIC:SpinCast1}{}
    \begin{enumerate}[label={\bf{}(\alph*)}]
    \item Use the diagram above to show that the curve must satisfy the
      differential equation
      $$
      \dfdx{y}{x}=\frac{\omega^2}{g} x
      $$
      and that the curve satisfying this must be a parabola.
    \item Using the value
      $g=9.8\sfrac{\text{meters}}{\text{second}^2}$, plot the parabola
      for $\omega{}=1, 2, 3$.  Do these plots coincide with your
      intuition about what should happen as the liquid rotates faster?
% What will happen to the parabola as $\omega$ increases? Do
%       you have any physical intuition about this problem? If so, is
%       your intuition consistent with the results of your analysis
%       above?
    \end{enumerate}
  \end{embeddedproblem-enumerate}

  Notice that for Problem~\#\ref{PIC:SpinCast1}, we simply gave you a
  formula for the centripetal force. We will derive this formula
  analytically later when we have extended the scope of our
  differentiation rules.

  Recall that by using a dynamical approach to constructing tangents
  Gilles Personne de Roberval was able to demonstrate that any light
  ray coming in parallel to the axis of a parabola will reflect to a
  single point (the focus of the parabola).  This is why the primary
  mirror in a telescope must be parabolic.

  In his derivation, Roberval used the focus-directrix definition of a
  parabola, rather than the equation we obtained in
  Problem~\#\ref{PIC:SpinCast1}.  That these two formulations are
  equivalent is not obvious.

  The next problem explores this fact.

    \pagebreak[4]
  \begin{embeddedproblem}{}
  \vskip-7mm{}
  \begin{wrapfigure}[7]{r}{2.5in}
    \vskip-4mm{}
    \captionsetup{labelformat=empty}
    \centerline{\includegraphics*[height=1.6in,width=2.5in]{../Figures/RobervalProb1}}
    \label{fig:RobervalProb1}
  \end{wrapfigure}
  Suppose that
  the point $X=(x,y)$ lies on the parabola with focus $F=(0,p)$ and
  directrix $y=-p$ as in the diagram above at the right.
    % \\
    % \centerline{\includegraphics*[height=2in,width=4in]{../Figures/RobervalProb1}}
    
    \begin{enumerate}[label={\bf{}(\alph*)}]{}
    \item Show that $X$ must satisfy the equation: $x^2=4py.$
    \item In Problem~\#\ref{PIC:SpinCast1} what should the angular
      velocity $\omega$ be to produce a focal length\aside{Having
        the correct focal length is of paramount importance in the
        performance of the telescope.  To give you a notion of how precise
        this must be, it took $3$ months to polish the surface of one of
        the primary mirrors in the Large Binocular Telescope to within an
        accuracy of $30$ nanometers ($3,000$ times thinner than a human
        hair).} of $p$?  
    \end{enumerate}
  \end{embeddedproblem}


  \begin{wrapfigure}[5]{r}{1.6in}
    \vskip-9mm{}
\captionsetup{labelformat=empty}
\centerline{\includegraphics*[height=1in,width=1.6in]{../Figures/TanSlope}}
\label{fig:}
 \end{wrapfigure}
  In antiquity the reflective properties of the conic sections
  (parabola, ellipse and hyperbola) were worked out via \emph{static}
  geometric arguments which are too laborious to give here. Following
  Roberval's lead we used \emph{dynamic}, but still geometric,
  arguments to accomplish the same task in
  Section~\ref{sec:roberv-conic-sect}.

  We now have everything we need to work them out again using
  Calculus. You can decide for yourself which method you find easier
  to work with\aside{This is not to suggest that you must, or even
    should, prefer the Calculus based method. Calculus, like Algebra
    or Geometry, is merely a tool. You should be proficient with all
    of your tools so that  you can use the tool
    that best fits the task at hand.}.

To apply calculus to the general problem, we will first need to find
  a formula for the angle between two lines in terms of their slopes.
  Have you noticed that the slope of a line is related to the tangent
  of the angle of elevation of the line?  This is readily seen in the
  diagram at the right.
  
  \begin{embeddedproblem}
      \label{problem:DiffOfTangents}
     \begin{wrapfigure}[]{r}{1.7in}
       \vskip-8mm{}
\captionsetup{labelformat=empty}
  \centerline{\includegraphics*[height=1.3in,width=1.7in]{../Figures/AngleBetweenLines}}
\label{fig:AngleBetweenLines}
\end{wrapfigure}
\vskip-8mm{}
Use the trigonometric identity\aside{We will see where this formula
  comes from in Problem~\#\ref{problem:DiffFormTan}.},
$$\tan(A-B)=\frac{\tan(A)-\tan(B)}{1+\tan(A)\tan(B)}$$ and the diagram
  at the right to show that
  $$  \tan(\theta) = \frac{m_2-m_1}{1+m_1m_2}.$$
\end{embeddedproblem}
\pagebreak[4]

   \begin{embeddedproblem}
     Consider the line tangent to the parabola $y=\frac{1}{4p}x^2$ at the point $(x_0,y_0)$.\\
     \centerline{\includegraphics*[height=1.4in,width=2.1in]{../Figures/ParabolaReflect}}
     Show that $\tan(\alpha)=\frac{2p}{x_0}=\tan(\beta)$ and use this
     to deduce the reflective property of a parabola.
   \end{embeddedproblem}

   The reflective property of the ellipse can be shown in a similar
   manner. Recall that an ellipse is the set of points in the plane,
   the sum of whose distances to two fixed points (foci) is a
   constant. Consider an ellipse with foci at the points $(-c,0)$ and
   $(c,0)$, and with constant sum $2a$ as seen in the next problem. Assume also that
   $a\gt c\gt 0$. (It is probably not clear why  this condition is
   necessary. See if you can figure it out as you work through this
   problem.)


   \begin{embeddedproblem}
     Use the diagram in the next problem and the fact that if $P$ lies
     on the ellipse, then $F_1P+PF_2=2a$ to show that $P$ must satisfy
     the equation\\
     \centerline{$\displaystyle
     \frac{x^2}{a^2}+\frac{y^2}{a^2-c^2}=1.
     $}
   \vskip-4mm{}
   \end{embeddedproblem}
   \vskip-4mm{}
   \begin{embeddedproblem}
        \vskip-10mm{}
     \centerline{\includegraphics*[height=1.5in,width=3.6in]{../Figures/EllipseReflect}}
     Given the tangent line to the ellipse
     $\frac{x^2}{a^2}+\frac{y^2}{a^2-c^2}=1$, an arbitrary point
     $(x_0, y_0)$ on the ellipse and the angles $\alpha$ and $\beta$
     as in the diagram above show that
     $$
     \tan(\alpha)=\frac{a^2-c^2}{cy_0}=\tan(\beta),
     $$
and use this to deduce the reflective property of an ellipse.
\end{embeddedproblem}

 The same calculations will work for the  hyperbola, but we'll take a
slightly different tack. It turns out that if one has a hyperbola and
an ellipse which have the same foci\aside{They are said to be  term
  \term{ confocal}.},
then wherever they intersect, they must be perpendicular. This is in
fact true for any set of confocal conic sections.

But before we get to that, we will need to derive the equation of a
hyperbola. This is similar to the derivation of the equation
for the ellipse.

\begin{embeddedproblem}
  \centerline{\includegraphics*[height=1in,width=2.7in]{../Figures/HyperbolaReflect}}
  Consider the point $P$ lying on the hyperbola with foci $F_1$ and
  $F_2$ and constant difference $2b$. This time we assume that $b\le
  c$. (Think about it.)

  Use the fact that $\abs{F_1P-PF_2}=2b$ to show that the point $P$ on
  the hyperbola must satisfy
  $$
  \frac{x^2}{b^2}-\frac{y^2}{c^2-b^2}=1.
  $$
\end{embeddedproblem}

\begin{embeddedproblem}
  Consider the point of intersection $P(x_0,y_0)$ of the following
  confocal ellipse and hyperbola seen in the diagram below.\\
  % \centerline{\includegraphics*[height=2in,width=3.2in]{../Figures/HypEllipConfocal}}
  \centerline{\includegraphics*[height=1.7in,width=2.7in]{../Figures/HypEllipConfocal}}
  Their equations, respectively, are
  $$
  \frac{x^2}{a^2}+\frac{y^2}{a^2-c^2}=1 \text{ and }
  \frac{x^2}{b^2}+\frac{y^2}{c^2-a^2}=1,
  $$
where this time $b\lt{}c\lt{} c$. (Why?)
       \begin{enumerate}[label={\bf{}(\alph*)}]{}
       \item Show that the tangent lines to these two curves at
         $P$ are perpendicular.
       \item Use this and the reflective property of an ellipse to
         show that any light ray directed toward one focus of the
         hyperbola will reflect off a branch of the hyperbola toward
         the other focus. 
            \end{enumerate}
\end{embeddedproblem}



%   To apply Calculus to this problem it is only necessary to observe
%   that when a light ray reflects from a point on a curve it reflects
%   in exactly the same fashion as if it were reflecting from the line
%   tangent to the curve at the same point. And, as we observed in
%   Section~\ref{sec:laziness-nature} when light reflects from a flat
%   mirror the angle of incidence will be equal to the angle of
%   reflection.

%   Before we address the conic sections it will be helpful to review
%   how to find the angle between two straight lines. 
  
%   \begin{ProficiencyDrill}
%     \label{drill:angle-between-lines}
%     This problem does not require Calculus.
%     \begin{enumerate}[label={\bf{}(\alph*)}]{}
%     \item Suppose that $m_1\ge m_2$. Show that the angle, $\theta$, between the
%       lines $y=m_1x$ and $y=m_2x$ is
%       $$\theta=\inverse\tan(m_1)-\inverse\tan(m_2).$$

%       What happens if
%       $m_2\ge m_1$?
%     \item Now show that the save formula works if $y=m_1x+b_1$ and
%       $y=m_2x+b_2$. That is, if the lines do not cross at the origin.
%     \end{enumerate}
%   \end{ProficiencyDrill}

    
%      The next problem asks you to verify the reflective properties of
%      the conic sections that we stated in
%      Section~\ref{sec:roberv-conic-sect}.
     
%     \begin{embeddedproblem-enumerate}[The Reflective Properties of the
%       Conic Sections]
%        \begin{enumerate}[label={\bf{}(\alph*)}]{}
%        \item {\bf[The Parabola]}
%          The graph of   $x^2-4py=0$ is a
%          parabola with its focus at $(0,p)$. In the diagram below the
%          point $(\alpha,\beta)$ is an arbitrary point on the
%          graph. Show that the angles $\phi$ and $\theta$ are equal.\\
%          \centerline{\includegraphics*[height=1.1in,width=2in]{../Figures/ParabolaReflect}}
%        \item {\bf{}[The Ellipse]} The graph of  $\frac{x^2}{a^2} +
%          \frac{y^2}{a^2-c^2}=1$ is an ellipse with its two foci at
%          $(-a,0)$ and $(a,0)$. In the diagram below the
%          point $(\alpha,\beta)$ is an arbitrary point on the
%          graph. Show that the angles $\phi$ and $\theta$ are equal.\\
%          \centerline{\includegraphics*[height=1.1in,width=2.5in]{../Figures/EllipseReflect}}
%      \item {\bf{}[The Hyperbola]} The graph of
%        $\frac{x^2}{a^2} - \frac{y^2}{c^2-a^2}=1$ is an hyperbola with
%        its two foci at $(-c,0)$ and $(c,0)$. In the diagram below the
%        point $(\alpha,\beta)$ is an arbitrary point on the
%        graph. Show that the angles $\phi$ and $\theta$ are equal.\\
%          \centerline{\includegraphics*[height=1.5in,width=2in]{../Figures/HyperbolaReflect}}
%        \end{enumerate}
% \end{embeddedproblem-enumerate}



  %%% Local Variables: 
  %%% outline-minor-mode: t
  %%% eval:(flyspell-mode)
  %%% TeX-master: "DiffCalc"
  %%% End: 



